% =========================================================
% Final lecture notes -- Time Series Analysis
% 6 chapters: Basic statistics, Decomposition, Stationarity,
%             Non-stationarity, ARMA, ARCH/GARCH
% Author : Aymen Ben Brik (aymen.benbrik@esprit.tn)
% Esprit School of Business
%
% Compile: pdflatex polycopie.tex (twice for TOC and cross-references)
% =========================================================
\documentclass[a4paper,11pt,oneside]{book}

% =========================================================
% Lecture notes preamble — Time Series Analysis module
% book class + tcolorbox + fancyhdr
% Author : Aymen Ben Brik (aymen.benbrik@esprit.tn)
% Esprit School of Business
% =========================================================

\usepackage[utf8]{inputenc}
\usepackage[T1]{fontenc}
\usepackage[english]{babel}
\usepackage{lmodern}
\usepackage{amsmath, amssymb, amsthm, mathtools, bm}
\usepackage{graphicx}
\usepackage{booktabs}
\usepackage{array}
\usepackage{multirow}
\usepackage{colortbl}
\usepackage{xcolor}
\usepackage{tikz}
\usetikzlibrary{positioning, arrows.meta, calc, shapes, fit, backgrounds, matrix, decorations.pathreplacing}
\usepackage{listings}
\usepackage[most]{tcolorbox}
\usepackage{geometry}
\geometry{a4paper, margin=2.5cm, headheight=14pt}
\usepackage{fancyhdr}
\usepackage[hidelinks]{hyperref}
\usepackage{microtype}
\usepackage{enumitem}

% --- Colour palette ---
\definecolor{espritBlue}{RGB}{30, 60, 110}
\definecolor{espritAccent}{RGB}{200, 80, 30}
\definecolor{boxDef}{RGB}{30, 60, 110}
\definecolor{boxTh}{RGB}{120, 30, 30}
\definecolor{boxProp}{RGB}{30, 100, 60}
\definecolor{boxEx}{RGB}{180, 120, 0}
\definecolor{boxRem}{RGB}{90, 90, 90}
\definecolor{boxWarn}{RGB}{200, 80, 30}

% --- Code listings ---
\definecolor{codebg}{RGB}{248,248,248}
\definecolor{codeframe}{RGB}{220,220,220}
\definecolor{keyword}{RGB}{0,82,155}
\definecolor{comment}{RGB}{88,110,117}
\definecolor{string}{RGB}{133,153,0}

\lstdefinestyle{pythonstyle}{
  language=Python,
  basicstyle=\ttfamily\small,
  keywordstyle=\color{keyword}\bfseries,
  commentstyle=\color{comment}\itshape,
  stringstyle=\color{string},
  backgroundcolor=\color{codebg},
  frame=single,
  rulecolor=\color{codeframe},
  frameround=tttt,
  breaklines=true,
  showstringspaces=false,
  tabsize=2
}
\lstdefinestyle{Rstyle}{
  language=R,
  basicstyle=\ttfamily\small,
  keywordstyle=\color{keyword}\bfseries,
  commentstyle=\color{comment}\itshape,
  stringstyle=\color{string},
  backgroundcolor=\color{codebg},
  frame=single,
  rulecolor=\color{codeframe},
  frameround=tttt,
  breaklines=true,
  showstringspaces=false,
  tabsize=2
}
\lstset{style=pythonstyle}

% --- Common macros ---
% =========================================================
% Common math macros — Time Series Analysis module
% Author : Aymen Ben Brik (aymen.benbrik@esprit.tn)
% Esprit School of Business
% =========================================================

% --- Standard sets ---
\providecommand{\R}{\mathbb{R}}
\providecommand{\N}{\mathbb{N}}
\providecommand{\Z}{\mathbb{Z}}
\providecommand{\Q}{\mathbb{Q}}
\providecommand{\C}{\mathbb{C}}

% --- Probability / statistics ---
\providecommand{\E}{\mathbb{E}}
\providecommand{\Prob}{\mathbb{P}}
\providecommand{\Var}{\mathrm{Var}}
\providecommand{\Cov}{\mathrm{Cov}}
\providecommand{\Corr}{\mathrm{Corr}}
\providecommand{\indep}{\perp\!\!\!\perp}
\providecommand{\iid}{\overset{\text{i.i.d.}}{\sim}}
\providecommand{\given}{\,\big\vert\,}

% --- Estimators / hat ---
\providecommand{\hatm}{\widehat{m}}
\providecommand{\hatsig}{\widehat{\sigma}}
\providecommand{\hatrho}{\widehat{\rho}}
\providecommand{\hatphi}{\widehat{\phi}}
\providecommand{\hattheta}{\widehat{\theta}}

% --- Time series notation ---
\providecommand{\Lop}{L}                      % lag operator
\providecommand{\Diff}{\Delta}                % difference operator
\providecommand{\AR}[1]{\mathrm{AR}(#1)}
\providecommand{\MAm}[1]{\mathrm{MA}(#1)}
\providecommand{\ARMA}[2]{\mathrm{ARMA}(#1,#2)}
\providecommand{\ARIMA}[3]{\mathrm{ARIMA}(#1,#2,#3)}
\providecommand{\SARIMA}[6]{\mathrm{SARIMA}(#1,#2,#3)(#4,#5,#6)}
\providecommand{\ARCH}[1]{\mathrm{ARCH}(#1)}
\providecommand{\GARCH}[2]{\mathrm{GARCH}(#1,#2)}
\providecommand{\WN}{\mathrm{WN}}              % white noise
\providecommand{\MSE}{\mathrm{MSE}}
\providecommand{\MAE}{\mathrm{MAE}}
\providecommand{\RMSE}{\mathrm{RMSE}}
\providecommand{\AIC}{\mathrm{AIC}}
\providecommand{\BIC}{\mathrm{BIC}}
\providecommand{\ACF}{\mathrm{ACF}}
\providecommand{\PACF}{\mathrm{PACF}}

% --- Linear algebra ---
\providecommand{\transp}[1]{#1^{\!\top}}
\providecommand{\norm}[1]{\left\lVert #1 \right\rVert}
\providecommand{\abs}[1]{\left\lvert #1 \right\rvert}
\providecommand{\inner}[2]{\langle #1, #2 \rangle}
\providecommand{\diag}{\mathrm{diag}}
\providecommand{\tr}{\mathrm{tr}}

% --- Bold vectors / matrices ---
\providecommand{\vx}{\bm{x}}
\providecommand{\vy}{\bm{y}}
\providecommand{\vz}{\bm{z}}
\providecommand{\vh}{\bm{h}}
\providecommand{\vu}{\bm{u}}
\providecommand{\vv}{\bm{v}}
\providecommand{\vw}{\bm{w}}
\providecommand{\vb}{\bm{b}}
\providecommand{\vW}{\bm{W}}
\providecommand{\vSigma}{\bm{\Sigma}}
\providecommand{\veps}{\bm{\varepsilon}}

% --- Author identity (reused on titlepages) ---
\providecommand{\auteurNom}{Aymen Ben Brik}
\providecommand{\auteurEmail}{aymen.benbrik@esprit.tn}
\providecommand{\auteurInstitution}{Esprit School of Business}
\providecommand{\auteurRole}{Module head}
\providecommand{\moduleNom}{Time Series Analysis}
\providecommand{\anneeUniversitaire}{2025--2026}


% --- Counters ---
\newcounter{boxcounter}[chapter]
\renewcommand{\theboxcounter}{\thechapter.\arabic{boxcounter}}

% --- tcolorbox styles (English labels) ---
\newtcolorbox[use counter=boxcounter]{defbox}[1][]{
  enhanced, breakable,
  colback=boxDef!5, colframe=boxDef, fonttitle=\bfseries,
  title={Definition~\theboxcounter\ifstrempty{#1}{}{ -- #1}},
  arc=2pt, boxrule=0.6pt
}
\newtcolorbox[use counter=boxcounter]{thbox}[1][]{
  enhanced, breakable,
  colback=boxTh!5, colframe=boxTh, fonttitle=\bfseries,
  title={Theorem~\theboxcounter\ifstrempty{#1}{}{ -- #1}},
  arc=2pt, boxrule=0.6pt
}
\newtcolorbox[use counter=boxcounter]{propbox}[1][]{
  enhanced, breakable,
  colback=boxProp!5, colframe=boxProp, fonttitle=\bfseries,
  title={Proposition~\theboxcounter\ifstrempty{#1}{}{ -- #1}},
  arc=2pt, boxrule=0.6pt
}
\newtcolorbox[use counter=boxcounter]{exbox}[1][]{
  enhanced, breakable,
  colback=boxEx!5, colframe=boxEx, fonttitle=\bfseries,
  title={Example~\theboxcounter\ifstrempty{#1}{}{ -- #1}},
  arc=2pt, boxrule=0.6pt
}
\newtcolorbox{rembox}[1][]{
  enhanced, breakable,
  colback=boxRem!5, colframe=boxRem, fonttitle=\bfseries,
  title={Remark\ifstrempty{#1}{}{ -- #1}},
  arc=2pt, boxrule=0.6pt
}
\newtcolorbox{warnbox}[1][]{
  enhanced, breakable,
  colback=boxWarn!5, colframe=boxWarn, fonttitle=\bfseries,
  title={Warning\ifstrempty{#1}{}{ -- #1}},
  arc=2pt, boxrule=0.6pt
}
\newtcolorbox{infobox}[1][]{
  enhanced, breakable,
  colback=espritBlue!5, colframe=espritBlue,
  arc=2pt, boxrule=0.5pt, left=4pt, right=4pt, top=3pt, bottom=3pt, #1
}

% --- Proof environment ---
\newtheoremstyle{proofstyle}{6pt}{6pt}{\itshape}{}{\bfseries}{.}{.5em}{}
\theoremstyle{proofstyle}

% --- Headers / footers ---
\pagestyle{fancy}
\fancyhf{}
\fancyhead[L]{\textsc{\moduleNom}}
\fancyhead[R]{\auteurNom}
\fancyfoot[C]{\thepage}
\fancyfoot[L]{\auteurInstitution}
\fancyfoot[R]{\auteurEmail}
\renewcommand{\headrulewidth}{0.4pt}
\renewcommand{\footrulewidth}{0.2pt}


\title{\moduleNom}
\author{\auteurNom}
\date{\anneeUniversitaire}

\begin{document}

% =========================================================
% Frontmatter: title page, foreword, table of contents
% =========================================================
\frontmatter
\pagestyle{empty}

% --- Title page ---
\begin{titlepage}
\centering
\vspace*{1.5cm}
{\scshape\large \auteurInstitution \par}
\vspace{0.3cm}
{\large Academic year \anneeUniversitaire \par}
\vspace{2.0cm}

\rule{\textwidth}{1.2pt}\par
\vspace{0.5cm}
{\Huge\bfseries \moduleNom \par}
\vspace{0.4cm}
{\Large Statistics $\cdot$ Decomposition $\cdot$ Stationarity $\cdot$
        Non-stationarity $\cdot$ ARMA $\cdot$ ARCH / GARCH \par}
\vspace{0.5cm}
\rule{\textwidth}{1.2pt}\par

\vspace{2.0cm}
{\large Lecture notes \par}

\vfill
{\Large\bfseries \auteurNom \par}
\vspace{0.4cm}
{\large \auteurRole \par}
\vspace{0.3cm}
{\large \href{mailto:\auteurEmail}{\auteurEmail} \par}

\vspace{1.0cm}
\end{titlepage}

% --- Copyright / version page ---
\thispagestyle{empty}
\vspace*{\fill}
\noindent\small
\textbf{\moduleNom} --- Lecture notes.\\[2pt]
Author: \auteurNom\ (\auteurRole), \auteurInstitution.\\[2pt]
Contact: \href{mailto:\auteurEmail}{\auteurEmail}.\\[2pt]
Academic year: \anneeUniversitaire.\\[6pt]
Compiled on: \today.\\[6pt]
Document distributed to students of the GAMMA, MDSI, MKD, MBA, LFG
and LFIG programmes of \auteurInstitution. Any reproduction or
public distribution requires the author's prior agreement.
\vspace*{1cm}
\newpage

% --- Foreword ---
\chapter*{Foreword}
\addcontentsline{toc}{chapter}{Foreword}
\pagestyle{fancy}

These lecture notes accompany the \emph{\moduleNom} course taught at
\auteurInstitution. They gather, in a single document, six
chapters that build progressively the modern toolbox of applied
time-series analysis.

\medskip
\textbf{Chapter~\ref{ch:basic-stats}} reviews the basic
statistical concepts (moments, estimation, sampling distributions,
inference) needed throughout the course. It serves both as a
refresher and as the conceptual reference for the chapters that
follow.

\medskip
\textbf{Chapter~\ref{ch:decomposition}} presents the
\emph{classical decomposition} of a time series into trend,
seasonality, and residual. Trend estimation by polynomial fitting
and centred moving averages, and seasonality estimation by
seasonal averaging or Fourier coefficients, are derived
explicitly.

\medskip
\textbf{Chapter~\ref{ch:stationarity}} introduces \emph{strict} and
\emph{weak} stationarity, the autocovariance and partial
autocorrelation functions, the Wold decomposition, and the
Box--Jenkins reading of the sample ACF/PACF. The Ljung--Box test
for white noise is stated and proved at the level required for
the labs.

\medskip
\textbf{Chapter~\ref{ch:nonstationarity}} addresses
\emph{non-stationary} series. The backshift $B$ and difference
$\Delta$ operators are introduced, integration order $I(d)$ is
defined, and the unit-root tests --- Augmented
Dickey--Fuller (ADF), KPSS, Phillips--Perron --- are presented in
detail. ARIMA$(p, d, q)$ and SARIMA$(p, d, q)(P, D, Q)_s$ models
close the chapter.

\medskip
\textbf{Chapter~\ref{ch:arma}} is the central chapter on ARMA
models. AR$(p)$, MA$(q)$ and ARMA$(p, q)$ are defined in lag-operator
form; stationarity and invertibility are stated as polynomial-root
conditions; and the full Box--Jenkins workflow ---
identification, MLE estimation, AIC/BIC selection, residual
diagnostics, forecasting --- is laid out.

\medskip
\textbf{Chapter~\ref{ch:arch-garch}} closes the module with
ARCH$(q)$ and GARCH$(p, q)$ models for conditional
heteroskedasticity. Volatility clustering, Engle's LM test,
quasi-MLE estimation, and one-day Value-at-Risk are presented;
the chapter mentions the EGARCH and GJR-GARCH extensions for the
leverage effect.

\medskip
Each chapter is delivered with a Beamer Metropolis 16:9 slide
deck, a Python lab (statement, solution, student skeleton,
reference code), and a project consolidating the chapter's notions.
A final integrative project, a reflective activity, and the
official ECUE module sheet (in French) complete the module.

\medskip\medskip
\noindent
Happy reading.\par\medskip
\hfill\auteurNom\par
\hfill\textit{\auteurEmail}\par

\vspace*{2cm}

% --- Table of contents ---
\tableofcontents
\thispagestyle{fancy}

% =========================================================
% Mainmatter: the six chapters
% =========================================================
\mainmatter

% =========================================================
% Chapter 1 — Basic statistical concepts
% Time Series Analysis — Aymen Ben Brik (aymen.benbrik@esprit.tn)
% Esprit School of Business
%
% Chapter content (no preamble), to be \input by polycopie.tex or
% chapitre1-statistics/chapitre-standalone.tex.
% =========================================================

\chapter{Basic statistical concepts}
\label{ch:basic-stats}

\section*{Introduction}
\addcontentsline{toc}{section}{Introduction}

Before studying time series proper, this chapter reviews the
statistical concepts on which the rest of the course relies:
\emph{moments} of a distribution, classical \emph{point estimators},
\emph{Method of Moments} and \emph{Maximum Likelihood Estimation},
sampling distributions, and basic \emph{inference} (confidence
intervals and hypothesis tests).

\medskip
The convention throughout the chapter is $X_1, \dots, X_n$ are
\emph{independent and identically distributed} (i.i.d.) draws from a
distribution with density $f(x;\theta)$ parameterised by some unknown
$\theta$. We observe a sample and want to learn $\theta$. Time series
will later violate the i.i.d.\ assumption, but the basic estimators
introduced here remain the building blocks.

% =========================================================
\section{Moments of a distribution}
\label{sec:moments}

\begin{defbox}[Raw and central moments]
\label{def:moments}
Let $X$ be a real-valued random variable with density $f$ and mean
$\mu = \E[X]$.
\begin{itemize}
  \item The \emph{raw $\ell$-th moment} is
        $m'_\ell \;=\; \E[X^\ell] \;=\; \int x^\ell f(x)\, dx$.
  \item The \emph{central $\ell$-th moment} is
        $m_\ell \;=\; \E[(X - \mu)^\ell]$.
\end{itemize}
\end{defbox}

\paragraph{Useful quantities.}
\begin{align*}
  \mu          &= m'_1, &
  \sigma^2     &= m_2 = \Var[X], \\[2pt]
  \text{skewness } S(X) &= \frac{m_3}{\sigma^3}, &
  \text{kurtosis } K(X) &= \frac{m_4}{\sigma^4}.
\end{align*}

\begin{rembox}[Interpretation]
\begin{itemize}
  \item $S(X)$ measures \emph{asymmetry}. $S = 0$ for any symmetric
        distribution; $S > 0$ for a right-skewed distribution.
  \item $K(X)$ measures \emph{tail-heaviness}. The Gaussian has $K = 3$;
        the \emph{excess kurtosis} $K - 3$ is therefore the standard
        comparison: $K - 3 > 0$ indicates heavier tails than the
        Gaussian (leptokurtic).
\end{itemize}
\end{rembox}

\begin{exbox}[Moments of the standard normal]
For $X \sim \mathcal{N}(\mu, \sigma^2)$ one shows by direct integration
or via the moment-generating function:
\[
  \E[X] = \mu, \quad \Var[X] = \sigma^2, \quad m_3 = 0, \quad m_4 = 3\,\sigma^4.
\]
Hence $S(X) = 0$ and $K(X) = 3$.
\end{exbox}

% =========================================================
\section{Classical estimators and sampling distributions}
\label{sec:estimators}

\subsection{Sample mean and sample variance}

\begin{defbox}[Empirical estimators]
Given i.i.d.\ data $X_1, \dots, X_n$:
\[
  \bar{X}_n \;=\; \frac{1}{n}\sum_{i=1}^{n} X_i,
  \qquad
  S_n^2 \;=\; \frac{1}{n-1}\sum_{i=1}^{n}(X_i - \bar{X}_n)^2.
\]
\end{defbox}

\begin{propbox}[Unbiasedness of $\bar{X}_n$ and $S_n^2$]
\label{prop:unbiased}
For i.i.d.\ data with finite second moments, $\E[\bar{X}_n] = \mu$ and
$\E[S_n^2] = \sigma^2$.
\end{propbox}

\begin{proof}
For the mean: $\E[\bar{X}_n] = \tfrac{1}{n}\sum_i \E[X_i] = \mu$ by
linearity. For the variance, expand:
\[
  \sum_i (X_i - \bar{X}_n)^2 = \sum_i X_i^2 - n\, \bar{X}_n^2.
\]
Take expectations: $\E[X_i^2] = \mu^2 + \sigma^2$ and
$\E[\bar{X}_n^2] = \mu^2 + \sigma^2 / n$. Hence
$\E[\sum_i(X_i - \bar{X}_n)^2] = n(\mu^2 + \sigma^2) - n(\mu^2 + \sigma^2/n)
= (n-1)\,\sigma^2$. The $n - 1$ denominator in $S_n^2$ ensures
$\E[S_n^2] = \sigma^2$.
\end{proof}

\begin{rembox}[Why $n - 1$?]
Replacing the unknown $\mu$ by $\bar{X}_n$ in the variance formula
``costs'' one degree of freedom. Dividing by $n$ would systematically
underestimate $\sigma^2$.
\end{rembox}

\subsection{Sampling distribution under normality}

\begin{thbox}[Cochran-type results]
\label{th:sampling-normal}
If $X_1, \dots, X_n \iid \mathcal{N}(\mu, \sigma^2)$, then:
\begin{enumerate}
  \item $\bar{X}_n \sim \mathcal{N}(\mu,\, \sigma^2 / n)$;
  \item $\dfrac{(n-1)\, S_n^2}{\sigma^2} \sim \chi^2_{n-1}$;
  \item $\bar{X}_n$ and $S_n^2$ are independent;
  \item $\dfrac{\bar{X}_n - \mu}{S_n / \sqrt{n}} \sim t_{n-1}$.
\end{enumerate}
\end{thbox}

\begin{rembox}
Result (4) follows from (1)--(3): the numerator is standard normal, the
denominator is the square root of an independent $\chi^2_{n-1}$ divided
by its degrees of freedom — the very definition of Student's $t$.
\end{rembox}

\subsection{Asymptotic results}

\begin{thbox}[Law of Large Numbers]
\label{th:lln}
For i.i.d.\ data with $\E[X_i] = \mu$:
\[
  \bar{X}_n \;\xrightarrow{p}\; \mu \quad \text{as } n \to \infty.
\]
\end{thbox}

\begin{thbox}[Central Limit Theorem]
\label{th:clt}
For i.i.d.\ data with $\E[X_i] = \mu$ and $\Var[X_i] = \sigma^2 < \infty$:
\[
  \frac{\bar{X}_n - \mu}{\sigma / \sqrt{n}}
  \;\xrightarrow{d}\;
  \mathcal{N}(0, 1) \quad \text{as } n \to \infty.
\]
\end{thbox}

\begin{rembox}[Practical consequence]
Even when the $X_i$ are not Gaussian, the sample mean behaves
approximately as a normal random variable for large $n$. The rule of
thumb $n \geq 30$ is conservative; in practice the convergence is
faster for symmetric, light-tailed distributions and slower for skewed
or heavy-tailed ones.
\end{rembox}

% =========================================================
\section{Estimation methods}
\label{sec:estimation}

\subsection{Method of Moments (MoM)}

\begin{defbox}[Method of Moments]
\label{def:mom}
Let $\theta \in \R^k$. The MoM estimator $\widehat{\theta}_n^{\text{MoM}}$
solves the system of equations
\[
  \E_\theta[X^p] \;=\; \tfrac{1}{n}\sum_{i=1}^{n} X_i^p,
  \qquad p = 1, 2, \dots, k.
\]
That is, theoretical moments are equated to sample moments.
\end{defbox}

\begin{exbox}[$\mathcal{N}(\mu, \sigma^2)$ via MoM]
With $\theta = (\mu, \sigma^2)$ and $k = 2$:
\begin{align*}
  \mu        &= \bar{X}_n, \\
  \mu^2 + \sigma^2 &= \tfrac{1}{n}\sum_i X_i^2.
\end{align*}
Subtracting yields
$\widehat{\sigma}^2_{\text{MoM}} = \tfrac{1}{n}\sum(X_i - \bar{X}_n)^2$,
which differs from $S_n^2$ by a factor $n / (n - 1)$.
\end{exbox}

\subsection{Maximum Likelihood Estimation (MLE)}

\begin{defbox}[Likelihood and MLE]
For i.i.d.\ data with density $f(x; \theta)$:
\[
  L(\theta) \;=\; \prod_{i=1}^{n} f(X_i; \theta),
  \qquad
  \ell(\theta) = \log L(\theta).
\]
The MLE is $\widehat{\theta}_n^{\text{MLE}}
\;=\; \arg\max_\theta\, \ell(\theta)$.
\end{defbox}

\begin{exbox}[$\mathcal{N}(\mu, \sigma^2)$ via MLE]
\label{ex:mle-gaussian}
The log-likelihood is
\[
  \ell(\mu, \sigma^2)
  = -\tfrac{n}{2}\log(2\pi\sigma^2)
    - \tfrac{1}{2\sigma^2}\sum_{i=1}^{n}(X_i - \mu)^2.
\]
Setting $\partial \ell / \partial \mu = 0$ gives
$\widehat{\mu}_{\text{MLE}} = \bar{X}_n$. Then setting
$\partial \ell / \partial \sigma^2 = 0$ gives
$\widehat{\sigma}^2_{\text{MLE}}
= \tfrac{1}{n}\sum(X_i - \bar{X}_n)^2$. So MLE coincides with MoM here,
and both are biased for $\sigma^2$ (their expectation is
$\tfrac{n-1}{n}\sigma^2$).
\end{exbox}

\begin{thbox}[Asymptotic normality of the MLE]
\label{th:mle-asymp}
Under regularity conditions (identifiability, smoothness, finite Fisher
information $I(\theta_0) > 0$):
\[
  \sqrt{n}\,\bigl(\widehat{\theta}_n^{\text{MLE}} - \theta_0\bigr)
  \;\xrightarrow{d}\;
  \mathcal{N}\!\bigl(0,\, I(\theta_0)^{-1}\bigr).
\]
\end{thbox}

\begin{rembox}[MoM vs.\ MLE]
\begin{itemize}
  \item MoM is computationally simple — just match moments.
  \item MLE is \emph{asymptotically efficient}: its variance attains
        the Cramér--Rao lower bound $I(\theta_0)^{-1}/n$.
  \item Both estimators are consistent (under regularity).
  \item For ``nice'' families (Gaussian, exponential, Poisson), MoM
        and MLE often coincide or are very close.
\end{itemize}
\end{rembox}

% =========================================================
\section{Multivariate data}
\label{sec:multivariate}

\subsection{Joint, marginal, and conditional distributions}

\begin{defbox}[Joint and marginal density]
For two random variables $X, Y$ with joint density $f(x, y)$:
\[
  f_X(x) \;=\; \int f(x, y)\, dy,
  \qquad
  f_{Y \mid X}(y \mid x) \;=\; \frac{f(x, y)}{f_X(x)}.
\]
The chain rule reads $f(x, y) = f_{Y \mid X}(y \mid x)\, f_X(x)$.
\end{defbox}

\begin{defbox}[Independence]
$X$ and $Y$ are \emph{independent} iff $f(x, y) = f_X(x)\, f_Y(y)$ for
all $x, y$.
\end{defbox}

\subsection{Covariance and correlation}

\begin{defbox}[Covariance and correlation]
\[
  \Cov(X, Y) = \E\bigl[(X - \E X)(Y - \E Y)\bigr],
  \qquad
  \Corr(X, Y) = \frac{\Cov(X, Y)}{\sqrt{\Var X \Var Y}}.
\]
\end{defbox}

\begin{propbox}[Sample versions]
\[
  \widehat{\Cov}(X, Y) = \tfrac{1}{n - 1}\sum_{i=1}^{n}
    (X_i - \bar X)(Y_i - \bar Y),
  \qquad
  \widehat{\Corr}(X, Y) = \frac{\widehat{\Cov}(X, Y)}{\widehat{\sigma}_X\,\widehat{\sigma}_Y}.
\]
The sample correlation lies in $[-1, 1]$.
\end{propbox}

\begin{warnbox}[Correlation captures only linear dependence]
$\Corr(X, Y) = 0$ does not imply independence. Counter-example: let
$X$ be symmetric around $0$ and $Y = X^2$. Then $\E[XY] = \E[X^3] = 0
= \E[X]\E[Y]$, hence $\Cov(X, Y) = 0$, yet $Y$ is a deterministic
function of $X$.
\end{warnbox}

% =========================================================
\section{Statistical inference}
\label{sec:inference}

\subsection{Confidence intervals for the mean}

\begin{defbox}[Confidence interval]
A $(1 - \alpha)$-confidence interval for $\mu$ is a random interval
$[\widehat{L}_n, \widehat{U}_n]$ such that
$\Prob\bigl(\widehat{L}_n \leq \mu \leq \widehat{U}_n\bigr) = 1 - \alpha$
in repeated sampling.
\end{defbox}

\paragraph{Gaussian sample, known $\sigma$ ($z$-interval).}
Let $X_1, \dots, X_n \iid \mathcal{N}(\mu, \sigma^2)$ with $\sigma$
known. By Theorem~\ref{th:sampling-normal} (1):
\[
  \mathrm{CI}_{1-\alpha}(\mu)
  = \left[\,\bar{X}_n - z_{1-\alpha/2}\,\tfrac{\sigma}{\sqrt{n}},\;
            \bar{X}_n + z_{1-\alpha/2}\,\tfrac{\sigma}{\sqrt{n}}\,\right].
\]

\paragraph{Gaussian sample, unknown $\sigma$ ($t$-interval).}
Replace $\sigma$ by $S_n$ and use Student's $t$ (Theorem~\ref{th:sampling-normal} (4)):
\[
  \mathrm{CI}_{1-\alpha}(\mu)
  = \left[\,\bar{X}_n - t_{n-1, 1-\alpha/2}\,\tfrac{S_n}{\sqrt{n}},\;
            \bar{X}_n + t_{n-1, 1-\alpha/2}\,\tfrac{S_n}{\sqrt{n}}\,\right].
\]

\subsection{Hypothesis testing}

\begin{defbox}[Test on the mean]
Two-sided test:
$H_0: \mu = \mu_0$ vs.\ $H_1: \mu \neq \mu_0$. Test statistic
\[
  T \;=\; \frac{\bar{X}_n - \mu_0}{S_n / \sqrt{n}}
  \;\sim\; t_{n-1} \text{ under } H_0.
\]
Reject $H_0$ at level $\alpha$ if $|T| > t_{n-1, 1-\alpha/2}$, i.e.\ the
$p$-value $2 \cdot \Prob(T_{n-1} > |T|) < \alpha$.
\end{defbox}

\begin{rembox}[Errors]
\begin{itemize}
  \item Type I error (size of the test): $\alpha = \Prob(\text{reject } H_0 \mid H_0)$.
  \item Type II error: $\beta = \Prob(\text{not reject } H_0 \mid H_1)$.
  \item \emph{Power} of the test: $1 - \beta$, the probability of
        correctly rejecting $H_0$ when $H_1$ is true.
\end{itemize}
\end{rembox}

\begin{exbox}[Numerical illustration]
Suppose $n = 50$ Atlanta monthly average temperatures with sample
mean $\bar{X}_{50} = 61.2$ and sample standard deviation $S_{50} = 16.8$.
Test $H_0 : \mu = 60$ against $H_1 : \mu \neq 60$ at $\alpha = 0.05$.
\[
  T \;=\; \frac{61.2 - 60}{16.8 / \sqrt{50}}
       \;\approx\; 0.505,
  \qquad
  t_{49, 0.975} \;\approx\; 2.010.
\]
Since $|T| < 2.010$, we do \emph{not} reject $H_0$;
$p\text{-value} \approx 2 \cdot \Prob(t_{49} > 0.505) \approx 0.616$.
\medskip
The corresponding 95\% CI is
$[61.2 - 4.78,\, 61.2 + 4.78] = [56.4,\, 66.0]$, which contains
$\mu_0 = 60$, consistent with the test conclusion.
\end{exbox}

\begin{warnbox}[Common pitfalls]
\begin{itemize}
  \item ``Failing to reject $H_0$'' is \emph{not} ``proving $H_0$''.
  \item A small $p$-value does not imply a large effect — it only
        signals statistical significance, not practical importance.
  \item Multiple testing inflates Type I error: use Bonferroni or FDR
        corrections when running many tests.
\end{itemize}
\end{warnbox}

% =========================================================
\section*{Summary}
\addcontentsline{toc}{section}{Summary}

\begin{center}
\begin{tabular}{lll}
\toprule
\textbf{Quantity} & \textbf{Estimator} & \textbf{Sampling distribution (Gaussian data)} \\
\midrule
$\mu$         & $\bar{X}_n = \tfrac{1}{n}\sum X_i$
              & $\mathcal{N}(\mu, \sigma^2/n)$ \\
$\sigma^2$    & $S_n^2 = \tfrac{1}{n-1}\sum (X_i - \bar X_n)^2$
              & $(n-1)S_n^2/\sigma^2 \sim \chi^2_{n-1}$ \\
Pivot         & $T = (\bar{X}_n - \mu)/(S_n/\sqrt{n})$
              & $t_{n-1}$ \\
\bottomrule
\end{tabular}
\end{center}

\medskip
\noindent
\textbf{What to remember.}
\begin{itemize}
  \item Two estimation routes: MoM (match moments) and MLE
        (maximise likelihood). Both consistent; MLE asymptotically
        efficient.
  \item Under normality, the sampling distribution of $\bar X$ is
        exactly Gaussian; the CLT extends this asymptotically.
  \item Confidence intervals and hypothesis tests on the mean rely on
        Student's $t$ when $\sigma$ is unknown.
\end{itemize}


% =========================================================
% Chapter 2 — Time series decomposition
% Author : Aymen Ben Brik (aymen.benbrik@esprit.tn)
% =========================================================

\chapter{Time series decomposition}
\label{ch:decomposition}

\section*{Introduction}
\addcontentsline{toc}{section}{Introduction}

The first step in time-series analysis is to \emph{decompose} the
observed data into structural components: a slowly varying
\emph{trend} $T_t$, a periodic \emph{seasonality} $S_t$, and a
stationary residual $\varepsilon_t$. The trend and seasonality carry
the deterministic structure; the residual captures the stochastic
fluctuations and is the input to the modelling techniques of
Chapters~\ref{ch:stationarity}--\ref{ch:arch-garch}.

This chapter presents:
\begin{enumerate}
  \item the additive and multiplicative decomposition models;
  \item three trend-estimation strategies (moving average, parametric
        polynomial regression, non-parametric LOESS / splines);
  \item three seasonality-estimation strategies (seasonal averaging,
        dummy regression, cosine--sine Fourier models);
  \item residual diagnostics and the full decomposition workflow.
\end{enumerate}

% =========================================================
\section{The decomposition model}
\label{sec:decomp-model}

\begin{defbox}[Additive vs.\ multiplicative model]
\label{def:decomp}
For an observed time series $(Y_t)_{t=1}^{T}$ with seasonal period $d$:
\[
  \text{Additive:}\quad Y_t = T_t + S_t + \varepsilon_t,
  \qquad
  \text{Multiplicative:}\quad Y_t = T_t \cdot S_t \cdot \varepsilon_t.
\]
Here $T_t$ is a slowly varying \emph{trend}, $S_t$ is a periodic
\emph{seasonal component} of period $d$ (i.e.\ $S_{t + d} = S_t$), and
$\varepsilon_t$ is a zero-mean stationary residual (additive case) or
a positive multiplicative perturbation with $\E[\varepsilon_t] = 1$
(multiplicative case).
\end{defbox}

\begin{rembox}
The two models differ qualitatively:
\begin{itemize}
  \item In the additive model, the seasonal swing has a
        \emph{constant amplitude} regardless of the level of the series.
  \item In the multiplicative model, the swing \emph{grows
        proportionally to $T_t$}.
  \item Taking logarithms converts a multiplicative model into an
        additive one:
        $\log Y_t = \log T_t + \log S_t + \log \varepsilon_t$.
\end{itemize}
This log transform is standard for positive series with growing
amplitude (sales, population, GDP).
\end{rembox}

\paragraph{Identifiability.} Without further constraints the
decomposition is not unique (a constant can move from $T$ to $S$).
Two standard normalisations:
\begin{itemize}
  \item $\sum_{k = 1}^{d} S_k = 0$ in the additive case (seasonality
        sums to zero over a full cycle);
  \item $\sum_{k = 1}^{d} S_k = d$ in the multiplicative case
        (seasonality averages to one).
\end{itemize}

% =========================================================
\section{Trend estimation}
\label{sec:trend}

\subsection{Symmetric (centred) moving average}

\begin{defbox}[Centred moving average]
\label{def:ma}
For an odd window $2k + 1$ with $k \in \N$:
\[
  \widehat{T}_t \;=\; \frac{1}{2k + 1}\sum_{i=-k}^{k} Y_{t+i},
  \qquad t = k+1, \dots, T - k.
\]
Boundary points $t = 1, \dots, k$ and $t = T - k + 1, \dots, T$ are
not estimated.
\end{defbox}

\begin{propbox}[The centred MA preserves linear trends]
\label{prop:ma-linear}
If $Y_t = a + b\,t$ exactly, then $\widehat{T}_t = a + b\,t$ for all
interior $t$.
\end{propbox}

\begin{proof}
Substitute $Y_{t+i} = a + b(t + i)$:
\[
  \widehat{T}_t = \tfrac{1}{2k+1}\sum_{i = -k}^{k}\bigl(a + b(t+i)\bigr)
  = a + b\,t + \tfrac{b}{2k+1}\underbrace{\sum_{i=-k}^{k} i}_{= 0}
  = a + b\,t.
\]
\end{proof}

\begin{propbox}[The centred MA cancels period-$(2k+1)$ seasonality]
\label{prop:ma-season}
If $S_t$ has period $d = 2k + 1$ and $\sum_{k=1}^{d} S_k = 0$, then
$\sum_{i = -k}^{k} S_{t + i} = 0$ for all $t$, hence the MA filters out
$S_t$ exactly.
\end{propbox}

\begin{rembox}
For \emph{even} period $d$ (e.g.\ $d = 12$ for monthly data), one uses
a $2 \times d$ centred MA:
$\widehat{T}_t = \tfrac{1}{2d}\bigl(\tfrac{1}{2}Y_{t-d/2} + Y_{t-d/2+1}
+ \dots + Y_{t+d/2-1} + \tfrac{1}{2}Y_{t+d/2}\bigr)$.
This combines two consecutive symmetric MAs of length $d$ to yield a
window of length $d + 1$ with end-points down-weighted.
\end{rembox}

\begin{exbox}[Numerical illustration]
$Y = (10, 12, 15, 14, 16, 18, 20)$, $k = 1$ (window 3):
\[
  \widehat{T}_2 = \tfrac{37}{3} \approx 12.33, \quad
  \widehat{T}_3 = \tfrac{41}{3} \approx 13.67, \quad
  \widehat{T}_4 = \tfrac{45}{3} = 15.00,
\]
\[
  \widehat{T}_5 = 16.00, \quad
  \widehat{T}_6 = 18.00.
\]
The estimated trend follows the upward drift of the series; $t = 1$
and $t = 7$ are lost.
\end{exbox}

\subsection{Parametric polynomial regression}

\begin{defbox}[Polynomial trend]
$T_t = \beta_0 + \beta_1\, t + \beta_2\, t^2 + \dots + \beta_p\, t^p$
estimated by ordinary least squares (OLS) on $(t, Y_t)$.
\end{defbox}

\begin{rembox}
\begin{itemize}
  \item $p = 1$ : straight line. Use only when the eyeball plot is
        consistent with a linear pattern.
  \item $p = 2, 3$ : flexible enough for most economic / climate data
        on a moderate horizon.
  \item $p \geq 4$ : usually overfit and behave poorly outside the
        sample range. Prefer LOESS or splines.
\end{itemize}
Use AIC / BIC, cross-validation or out-of-sample MSE to select $p$.
\end{rembox}

\subsection{Non-parametric regression}

\begin{defbox}[LOESS / local polynomial regression]
At each $t$ fit a low-degree polynomial on a \emph{neighbourhood} of
$t$, with weights decreasing in $\abs{t' - t}$ (e.g.\ tricube kernel).
The fraction of data used in each local fit is the \emph{span}
parameter $\alpha \in (0, 1)$.
\end{defbox}

\paragraph{Why non-parametric?}
\begin{itemize}
  \item Captures arbitrary smooth trends without imposing a global
        functional form.
  \item Robust variants (\texttt{loess(\dots, family="symmetric")}) are
        less affected by outliers.
  \item Span $\alpha$ small $\Rightarrow$ wiggly fit; large $\Rightarrow$
        over-smooth. Default $\alpha = 0.75$ is often too aggressive
        for time-series trend extraction; $\alpha = 0.1$--$0.3$ is
        usually preferred.
\end{itemize}

\begin{rembox}[Other smoothers]
Cubic smoothing splines minimise
$\sum (Y_t - g(t))^2 + \lambda \int g''(s)^2\,ds$ over twice-differentiable
$g$. The penalty $\lambda$ trades off fidelity vs.\ smoothness;
$\lambda \to 0$ recovers interpolation, $\lambda \to \infty$ a linear
fit. Wavelets offer multi-resolution alternatives but are seldom needed
for a one-step decomposition.
\end{rembox}

% =========================================================
\section{Seasonality estimation}
\label{sec:seasonality}

\subsection{Seasonal averaging}

\begin{defbox}[Seasonal averaging after detrending]
\label{def:seas-avg}
Let $\widetilde{Y}_t = Y_t - \widehat{T}_t$ be the detrended series.
For each season $k = 1, \dots, d$ define
\[
  \widehat{S}_k \;=\; \frac{1}{n_k}\sum_{t : \text{season}(t) = k}
                    \widetilde{Y}_t,
\]
where $n_k$ is the number of observations falling in season $k$.
Re-centre by subtracting the overall mean to ensure $\sum_k \widehat{S}_k = 0$.
\end{defbox}

\begin{exbox}[Monthly seasonal index]
With monthly data ($d = 12$) over 10 years, $n_k = 10$ for each $k$.
The estimate $\widehat{S}_{\text{July}}$ is the average of the
ten July deviations from the trend.
\end{exbox}

\subsection{Dummy-variable regression}

\begin{defbox}[Seasonal dummies]
\label{def:dummies}
Let $D_{k,t} = \mathbf{1}_{\{\text{season}(t) = k\}}$ for $k = 2, \dots, d$.
Estimate by OLS:
\[
  Y_t = \beta_0 + \sum_{k=2}^{d} \beta_k\, D_{k,t} + \varepsilon_t.
\]
The omitted season ($k = 1$) is the \emph{baseline}; $\beta_k$
measures the seasonal deviation of season $k$ from the baseline.
\end{defbox}

\begin{rembox}[Why $d - 1$ dummies?]
Including all $d$ dummies plus an intercept produces perfect
multicollinearity ($\sum_k D_{k,t} = 1$ for all $t$). Drop one season
or drop the intercept.
\end{rembox}

\begin{rembox}[Combining trend and seasonality]
A standard ``one-shot'' specification estimates trend and seasonality
\emph{jointly}:
\[
  Y_t = \beta_0 + \beta_1\, t + \sum_{k=2}^{d} \gamma_k\, D_{k,t} + \varepsilon_t.
\]
This avoids the residual contamination that arises when one estimates
the trend first by MA then the seasonality from the detrended series.
\end{rembox}

\subsection{Cosine--sine (Fourier) model}

\begin{defbox}[Trigonometric seasonality]
\label{def:fourier}
With angular frequency $\omega_d = 2\pi / d$:
\[
  S_t = \sum_{j=1}^{J}\bigl[a_j \cos(j\,\omega_d\, t) + b_j \sin(j\,\omega_d\, t)\bigr],
  \qquad J \leq \lfloor d / 2 \rfloor.
\]
The coefficients $(a_j, b_j)$ are estimated by OLS in the regression
$Y_t = \beta_0 + S_t + \varepsilon_t$ (or jointly with a trend).
\end{defbox}

\begin{propbox}[Equivalence with dummies]
\label{prop:fourier-equiv}
With $J = \lfloor d/2 \rfloor$ harmonics (and the constant), the
trigonometric basis spans the same subspace as the seasonal dummies,
hence yields the same fitted seasonal pattern.
\end{propbox}

\begin{rembox}
The Fourier model is \emph{parsimonious}: with $J = 1$ harmonic only,
seasonal swings are described by 2 parameters instead of $d - 1$. This
helps when $d$ is large (e.g.\ $d = 365$ for daily data with annual
seasonality).
\end{rembox}

% =========================================================
\section{Residuals and diagnostics}
\label{sec:residuals}

After decomposition, the residual is
\[
  \widehat{\varepsilon}_t \;=\; Y_t - \widehat{T}_t - \widehat{S}_t.
\]

\begin{rembox}[What to check]
\begin{itemize}
  \item \textbf{No remaining trend.} Plot $\widehat{\varepsilon}_t$ vs.\ $t$;
        any systematic drift means under-fit trend.
  \item \textbf{No remaining seasonality.} Compute the autocorrelation
        function (Chapter 3) and look for spikes at lag $d$ (and its
        multiples).
  \item \textbf{Constant variance (homoscedasticity).} A funnel-shaped
        plot signals heteroscedasticity, suggesting a multiplicative
        rather than additive model — or ARCH/GARCH effects (Chapter 6).
  \item \textbf{Approximate normality.} Histogram + QQ-plot. Important
        for the validity of $t$-based confidence intervals.
\end{itemize}
\end{rembox}

\begin{warnbox}
A common mistake is to declare the decomposition successful as soon as
the residual ``looks small''. The residual must be \emph{structureless}
— that is, free of autocorrelation, of remaining seasonality and of
heteroscedasticity. The next chapters provide the formal tools to
check this.
\end{warnbox}

% =========================================================
\section{Decomposition workflow}
\label{sec:workflow}

The full pipeline:

\begin{enumerate}
  \item \textbf{Plot} $Y_t$ versus $t$. Decide additive vs.\ multiplicative
        from the visual amplitude pattern.
  \item If multiplicative, apply the log transform.
  \item \textbf{Estimate the trend} $\widehat{T}_t$:
        \begin{itemize}
          \item moving average (centred, $2 \times d$ for even $d$);
          \item polynomial regression ($p = 1, 2$ or $3$);
          \item LOESS / spline.
        \end{itemize}
  \item \textbf{Detrend:} $\widetilde{Y}_t = Y_t - \widehat{T}_t$.
  \item \textbf{Estimate the seasonality} $\widehat{S}_t$ from $\widetilde{Y}_t$:
        \begin{itemize}
          \item seasonal averaging (re-centred);
          \item dummy regression;
          \item cosine--sine model.
        \end{itemize}
  \item \textbf{Residual:}
        $\widehat{\varepsilon}_t = Y_t - \widehat{T}_t - \widehat{S}_t$.
  \item \textbf{Diagnostics} on $\widehat{\varepsilon}_t$ (Section~\ref{sec:residuals}).
  \item If diagnostics pass: model the residual with ARMA (Chapter 5).
\end{enumerate}

% =========================================================
\section*{Summary}
\addcontentsline{toc}{section}{Summary}

\begin{center}
\begin{tabular}{lll}
\toprule
\textbf{Component} & \textbf{Method} & \textbf{Pros / Cons} \\
\midrule
Trend       & Moving average ($2k+1$)      & simple, loses $k$ ends \\
            & Polynomial regression        & global, risk of overfit \\
            & LOESS / splines              & flexible, span tuning \\
\midrule
Seasonality & Seasonal averaging           & simple, no model \\
            & Dummy regression             & joint with trend \\
            & Cosine--sine (Fourier)       & parsimonious, multi-harmonic \\
\bottomrule
\end{tabular}
\end{center}

\medskip
\noindent\textbf{What to remember.}
\begin{itemize}
  \item Decomposition is the \emph{first step} of any time-series
        analysis; the rest of the course works on the residual.
  \item Choose additive or multiplicative based on the amplitude
        pattern of the seasonality.
  \item Estimate trend and seasonality jointly (dummy or Fourier
        regression) when feasible — it controls residual contamination.
\end{itemize}


% =========================================================
% Chapter 3 — Stationary time series
% Author : Aymen Ben Brik (aymen.benbrik@esprit.tn)
% =========================================================

\chapter{Stationary time series}
\label{ch:stationarity}

\section*{Introduction}
\addcontentsline{toc}{section}{Introduction}

After decomposing the observed series into trend, seasonality and
residual (Chapter~\ref{ch:decomposition}), we usually want the residual
to be \emph{stationary} so that the modelling tools of
Chapter~\ref{ch:arma} apply. This chapter formalises stationarity,
introduces the autocovariance and autocorrelation functions (ACF and
PACF), states the Wold decomposition theorem, and provides
\emph{white-noise tests} (Box--Pierce, Ljung--Box) to validate model
residuals.

% =========================================================
\section{Stationarity}
\label{sec:stationarity}

\subsection{Strict and weak stationarity}

\begin{defbox}[Strict (strong) stationarity]
\label{def:strict-stat}
A process $\{X_t\}_{t \in \Z}$ is \emph{strictly stationary} if, for
every integer $k$, every $t_1 < \dots < t_k$ and every $h \in \Z$,
\[
  (X_{t_1}, X_{t_2}, \dots, X_{t_k})
  \overset{\text{d}}{=}
  (X_{t_1 + h}, X_{t_2 + h}, \dots, X_{t_k + h}).
\]
\end{defbox}

\begin{defbox}[Weak (covariance) stationarity]
\label{def:weak-stat}
A process $\{X_t\}$ with $\E[X_t^2] < \infty$ is \emph{weakly
stationary} if there exist $\mu \in \R$ and a function
$\gamma : \Z \to \R$ such that
\[
  \E[X_t] = \mu, \qquad
  \Cov(X_t, X_{t+h}) = \gamma(h) \quad \text{for all } t, h \in \Z.
\]
The function $\gamma$ is the \emph{autocovariance function} (ACVF).
\end{defbox}

\begin{propbox}[Strict $\Rightarrow$ weak (under $L^2$)]
\label{prop:strict-weak}
If $\{X_t\}$ is strictly stationary with $\E[X_t^2] < \infty$, then
$\{X_t\}$ is weakly stationary.
\end{propbox}

\begin{proof}
By strict stationarity, $X_t \overset{\text{d}}{=} X_0$, so
$\E[X_t] = \E[X_0] = \mu$ does not depend on $t$. Similarly,
$(X_t, X_{t+h}) \overset{\text{d}}{=} (X_0, X_h)$, hence
$\Cov(X_t, X_{t+h}) = \Cov(X_0, X_h)$ depends only on $h$.
\end{proof}

\begin{rembox}[The converse is false in general]
A process can be weakly stationary without being strictly stationary
if it has time-varying higher-order moments. For Gaussian processes
(only mean and covariance matter), the two notions coincide.
\end{rembox}

\subsection{Two canonical examples}

\begin{defbox}[White noise]
\label{def:wn}
$(\varepsilon_t)$ is a \emph{white noise} $\WN(0, \sigma^2)$ if
$\E[\varepsilon_t] = 0$, $\Var[\varepsilon_t] = \sigma^2$ and
$\Cov(\varepsilon_t, \varepsilon_{t+h}) = 0$ for $h \neq 0$.
A \emph{strong white noise} requires the $\varepsilon_t$ to be
i.i.d.; a Gaussian white noise has $\varepsilon_t \iid \mathcal{N}(0, \sigma^2)$.
\end{defbox}

\begin{exbox}[Random walk is \emph{not} stationary]
\label{ex:random-walk}
Let $\varepsilon_t \iid \WN(0, \sigma^2)$ and
$X_t = X_{t-1} + \varepsilon_t$ with $X_0 = 0$. Then
$X_t = \sum_{s = 1}^{t} \varepsilon_s$ and
\[
  \E[X_t] = 0,
  \quad
  \Var[X_t] = t\, \sigma^2,
  \quad
  \Cov(X_t, X_{t+h}) = t\, \sigma^2 \text{ for } h \geq 0.
\]
The variance grows linearly in $t$, so $X_t$ is non-stationary. This
is the canonical example used by unit-root tests in
Chapter~\ref{ch:nonstationarity}.
\end{exbox}

% =========================================================
\section{Autocovariance and autocorrelation}
\label{sec:acf}

\begin{defbox}[ACVF and ACF]
For a weakly stationary process with mean $\mu$:
\[
  \gamma(h) = \Cov(X_t, X_{t+h}),
  \qquad
  \rho(h) = \frac{\gamma(h)}{\gamma(0)} \;\in\; [-1, 1].
\]
\end{defbox}

\begin{propbox}[Properties of $\gamma$]
\label{prop:acvf}
\begin{enumerate}
  \item $\gamma(0) = \Var[X_t] \geq 0$;
  \item $\abs{\gamma(h)} \leq \gamma(0)$ (Cauchy--Schwarz);
  \item $\gamma$ is even: $\gamma(-h) = \gamma(h)$;
  \item $\gamma$ is positive semidefinite: for any $n$ and any
        $a_1, \dots, a_n \in \R$,
        $\sum_{i, j = 1}^{n} a_i\, a_j\, \gamma(t_i - t_j) \geq 0$.
\end{enumerate}
\end{propbox}

\begin{proof}
(1) and (3) follow from the definition. (2) follows from Cauchy--Schwarz
applied to $X_t - \mu$ and $X_{t+h} - \mu$. (4) is the standard fact
that any matrix of the form $\Sigma_{ij} = \gamma(t_i - t_j)$ is the
covariance matrix of $(X_{t_1}, \dots, X_{t_n})$, hence positive
semidefinite.
\end{proof}

\subsection{Sample ACF and Bartlett bands}

\begin{defbox}[Sample ACVF and ACF]
For $X_1, \dots, X_T$ with sample mean $\bar X$:
\[
  \widehat\gamma(h) = \frac{1}{T}\sum_{t = 1}^{T - h}
    (X_t - \bar X)(X_{t+h} - \bar X),
  \qquad
  \widehat\rho(h) = \frac{\widehat\gamma(h)}{\widehat\gamma(0)}.
\]
\end{defbox}

\begin{thbox}[Bartlett's formula for white noise]
\label{th:bartlett}
If $\{X_t\}$ is a white noise with finite fourth moment, then for any
fixed $h \geq 1$:
\[
  \sqrt{T}\,\widehat\rho(h) \;\xrightarrow{d}\; \mathcal{N}(0, 1).
\]
The standard test bands $\pm 1.96 / \sqrt{T}$ on the ACF plot are
asymptotic 95\% confidence intervals \emph{under the white-noise null}.
\end{thbox}

\begin{rembox}
Bartlett's formula gives a more general result for non-white processes
(involves a sum over the squared autocorrelations); see Brockwell \&
Davis, \emph{Time Series: Theory and Methods}, Chapter 7.
\end{rembox}

% =========================================================
\section{Partial autocorrelation}
\label{sec:pacf}

\begin{defbox}[Partial autocorrelation function]
\label{def:pacf}
The PACF $\alpha(h)$ at lag $h$ is the correlation between $X_t$ and
$X_{t+h}$ \emph{after removing the linear effect} of the intermediate
observations $X_{t+1}, \dots, X_{t+h-1}$:
\[
  \alpha(h) = \Corr\bigl(X_t - \widehat X_t,\; X_{t+h} - \widehat X_{t+h}\bigr),
\]
where $\widehat X_t$ and $\widehat X_{t+h}$ are the linear projections
of $X_t$ and $X_{t+h}$ on the span of $\{X_{t+1}, \dots, X_{t+h-1}\}$.
\end{defbox}

\begin{rembox}[Recursive computation]
The PACF is computed by the \emph{Durbin--Levinson} algorithm: it
solves the Yule--Walker equations of order $h$ and reads off the
coefficient of $X_{t+h-h} = X_t$ in the linear regression of $X_{t+h}$
on $X_{t+1}, \dots, X_{t+h}$. Most software (Python
\texttt{statsmodels}, R \texttt{stats}) returns the PACF in $O(h)$
operations per lag.
\end{rembox}

% =========================================================
\section{ACF / PACF signatures of standard processes}
\label{sec:acf-signatures}

\begin{center}
\begin{tabular}{lll}
\toprule
\textbf{Process} & \textbf{$\rho(h)$ behaviour} & \textbf{$\alpha(h)$ behaviour} \\
\midrule
$\WN$           & $\rho(h) = 0$ for $h \neq 0$         & $\alpha(h) = 0$ for $h \geq 1$ \\
$\AR{p}$        & decays (exp.\ / sinusoidal)          & cuts off at $h = p$ \\
$\MAm{q}$       & cuts off at $h = q$                  & decays \\
$\ARMA{p}{q}$   & decays after $h > q$                 & decays after $h > p$ \\
\bottomrule
\end{tabular}
\end{center}

\medskip
\noindent\textbf{Box--Jenkins identification.} Plot the sample ACF and
PACF, then use the table above to read off candidate orders $(p, q)$.

\subsection{AR(1) and MA(1) — closed-form ACF}

\begin{exbox}[AR(1) ACF]
\label{ex:ar1-acf}
For $X_t = \phi X_{t-1} + \varepsilon_t$ with
$\varepsilon_t \sim \WN(0, \sigma^2)$ and $\abs\phi < 1$:
\[
  \rho(h) = \phi^{\abs h}, \quad h \in \Z,
  \qquad
  \alpha(1) = \phi, \quad \alpha(h) = 0 \text{ for } h \geq 2.
\]
The ACF decays exponentially; the PACF cuts off at lag 1.
\end{exbox}

\begin{exbox}[MA(1) ACF]
\label{ex:ma1-acf}
For $X_t = \varepsilon_t + \theta \varepsilon_{t-1}$ with
$\varepsilon_t \sim \WN(0, \sigma^2)$:
\[
  \rho(0) = 1, \qquad
  \rho(1) = \frac{\theta}{1 + \theta^2}, \qquad
  \rho(h) = 0 \text{ for } h \geq 2.
\]
The ACF cuts off at lag 1; the PACF decays.
\end{exbox}

\begin{proof}[Proof of $\rho(1) = \theta / (1 + \theta^2)$ in MA(1)]
$\Var[X_t] = \E[(\varepsilon_t + \theta\varepsilon_{t-1})^2]
= (1 + \theta^2)\sigma^2$. Then
$\Cov(X_t, X_{t+1}) = \E[(\varepsilon_t + \theta\varepsilon_{t-1})
(\varepsilon_{t+1} + \theta\varepsilon_t)] = \theta\,\sigma^2$ since
$\varepsilon_t$'s are uncorrelated. Hence $\rho(1) =
\theta\sigma^2 / [(1 + \theta^2)\sigma^2] = \theta / (1 + \theta^2)$.
\end{proof}

% =========================================================
\section{Wold decomposition}
\label{sec:wold}

\begin{thbox}[Wold's theorem, 1938]
\label{th:wold}
Every \emph{purely non-deterministic} weakly stationary process
$\{X_t\}$ admits a unique representation
\[
  X_t \;=\; \mu \;+\; \sum_{j = 0}^{\infty} \psi_j\, \varepsilon_{t - j},
  \qquad
  \psi_0 = 1, \qquad \sum_{j = 0}^{\infty} \psi_j^2 < \infty,
\]
where $(\varepsilon_t)$ is a white-noise process (the
\emph{innovations}) of variance
$\sigma^2 = \Var[\varepsilon_t]$.
\end{thbox}

\begin{rembox}[Why it matters]
\begin{itemize}
  \item \textbf{Universality:} every stationary process can be viewed
        as an \emph{infinite MA process}.
  \item \textbf{Justification of ARMA:} ARMA models are
        finite-dimensional approximations of the Wold representation
        (with a rational $\psi(z)$).
  \item \textbf{Forecasting:} optimal linear predictors and their
        prediction errors are obtained from the Wold coefficients.
\end{itemize}
\end{rembox}

% =========================================================
\section{White-noise tests}
\label{sec:wn-tests}

After fitting a decomposition or an ARMA model, the residuals
$\widehat\varepsilon_t$ should look like white noise. Two classical
tests check this jointly across several lags.

\begin{defbox}[Box--Pierce and Ljung--Box statistics]
\label{def:lb}
For sample ACF $\widehat\rho(h)$ of the residuals up to lag $m$:
\[
  Q_{\text{BP}} = T \sum_{h = 1}^{m} \widehat\rho(h)^2,
  \qquad
  Q_{\text{LB}} = T(T + 2) \sum_{h = 1}^{m} \frac{\widehat\rho(h)^2}{T - h}.
\]
\end{defbox}

\begin{thbox}[Asymptotic distribution under $H_0$]
\label{th:lb}
Under $H_0 :$ \emph{the residuals are white noise}, both statistics
are asymptotically $\chi^2_{m - p - q}$, where $p + q$ is the number
of fitted ARMA parameters (zero if no parametric model has been
fitted).
\end{thbox}

\paragraph{Decision rule.} Reject $H_0$ at level $\alpha$ if
$Q > \chi^2_{m - p - q,\, 1 - \alpha}$, equivalently $p$-value $< \alpha$.

\paragraph{Choosing $m$.} A few rules of thumb:
\begin{itemize}
  \item For yearly data: $m \in \{10, 20\}$.
  \item For monthly data: $m \in \{12, 24, 36\}$ (covers seasonality).
  \item For high-frequency data: $m$ proportional to expected
        autocorrelation length.
\end{itemize}

\begin{warnbox}
The Ljung--Box test is preferred over Box--Pierce in finite samples:
its small-sample distribution is closer to the $\chi^2$ asymptotic
limit, especially when $m / T$ is not negligible.
\end{warnbox}

% =========================================================
\section*{Summary}
\addcontentsline{toc}{section}{Summary}

\noindent\textbf{Key takeaways.}
\begin{itemize}
  \item \emph{Strict stationarity} = invariance of all finite-dimensional
        distributions; \emph{weak stationarity} = only first two moments.
        Strict $\Rightarrow$ weak (under $L^2$); converse holds for
        Gaussian processes.
  \item The \emph{autocovariance} $\gamma(h)$ describes the dependence
        structure; the \emph{autocorrelation} $\rho(h)$ is its
        normalised version.
  \item ACF and PACF have characteristic patterns that allow
        \emph{Box--Jenkins identification} of AR(p) / MA(q) / ARMA(p,q).
  \item \emph{Wold decomposition}: every stationary process is an
        infinite moving average of innovations.
  \item \emph{Ljung--Box test} validates that residuals look like
        white noise.
\end{itemize}


% =========================================================
% Chapter 4 — Non-stationary time series
% Author : Aymen Ben Brik (aymen.benbrik@esprit.tn)
% =========================================================

\chapter{Non-stationary time series}
\label{ch:nonstationarity}

\section*{Introduction}
\addcontentsline{toc}{section}{Introduction}

Most real-world time series are \emph{not} stationary: GDP grows,
prices drift, populations expand, atmospheric CO$_2$ accumulates.
Applying the tools of Chapter~\ref{ch:stationarity} (ACF, ARMA) to a
raw non-stationary series gives misleading results — the famous
\emph{spurious regression} phenomenon. This chapter formalises
non-stationarity, introduces the algebraic operators (backshift,
difference) and the unit-root tests (ADF, KPSS) needed to detect it,
and culminates in the \emph{ARIMA} and \emph{SARIMA} families that
combine differencing with the ARMA structure of
Chapter~\ref{ch:arma}.

% =========================================================
\section{Limits of the stationarity assumption}
\label{sec:non-stat-limits}

\begin{warnbox}[Spurious regression]
\label{warn:spurious}
If $X_t$ and $Y_t$ are two \emph{independent} random walks and one
fits the linear regression $Y_t = \beta_0 + \beta_1 X_t + u_t$ by
OLS:
\begin{itemize}
  \item the standard $t$-statistic for $\widehat\beta_1$ rejects the
        null $\beta_1 = 0$ at any level with probability tending to
        \emph{one} as the sample size grows;
  \item the $R^2$ stays large (often $> 0.5$);
\end{itemize}
even though the two series are \emph{independent}. Granger \& Newbold
(1974) called this the \emph{spurious regression} effect; it is
caused by the non-stationarity of $X$ and $Y$.
\end{warnbox}

\begin{rembox}[Practical consequence]
Before applying any classical inference (regression, ARMA, $t$-test),
\emph{check stationarity}. If the series is non-stationary, transform
it first (typically by differencing).
\end{rembox}

\subsection{Trend stationary vs.\ difference stationary}

\begin{defbox}[TS and DS processes]
\label{def:ts-ds}
\begin{itemize}
  \item \emph{Trend stationary} (TS):
        $X_t = f(t) + \varepsilon_t$ where $f$ is deterministic and
        $\varepsilon_t$ is stationary. Shocks have \emph{transient}
        effects.
  \item \emph{Difference stationary} (DS):
        $\Delta X_t$ is stationary, but $X_t$ is not. The series has
        a \emph{unit root}; shocks have \emph{permanent} effects.
\end{itemize}
\end{defbox}

\begin{rembox}
TS series are made stationary by \emph{detrending}; DS series, by
\emph{differencing}. Visual inspection alone cannot reliably
distinguish the two — a slowly trending DS process looks identical to
a TS one. We use \emph{unit-root tests} (Section~\ref{sec:unit-roots}).
\end{rembox}

% =========================================================
\section{Backshift and difference operators}
\label{sec:operators}

\begin{defbox}[Backshift operator]
\label{def:backshift}
$B X_t = X_{t-1}$, and $B^k X_t = X_{t - k}$.
For polynomials $\phi(z) = 1 - \phi_1 z - \dots - \phi_p z^p$ and
$\theta(z) = 1 + \theta_1 z + \dots + \theta_q z^q$, the AR(p) and
MA(q) processes are written
$\phi(B) X_t = \varepsilon_t$ and $X_t = \theta(B) \varepsilon_t$.
\end{defbox}

\begin{propbox}[Stationarity through the AR polynomial]
\label{prop:ar-stationary}
The AR(p) process $\phi(B) X_t = \varepsilon_t$ is (weakly) stationary
if and only if every root of $\phi(z) = 0$ lies \emph{outside} the
closed unit disk.
\end{propbox}

\begin{exbox}[AR(1)]
$\phi(z) = 1 - \phi z$ has root $z = 1 / \phi$.
$\abs{1 / \phi} > 1 \Leftrightarrow \abs\phi < 1$. We recover the
familiar AR(1) stationarity condition.
\end{exbox}

\begin{defbox}[Difference operator]
\label{def:difference}
$\Delta = 1 - B$, so $\Delta X_t = X_t - X_{t-1}$.
$\Delta^d = (1 - B)^d$ is the $d$-th difference.
The seasonal difference of period $s$ is
$\Delta_s = 1 - B^s$.
\end{defbox}

\begin{exbox}[Higher-order differencing]
$\Delta^2 X_t = (1 - B)^2 X_t = X_t - 2 X_{t-1} + X_{t-2}$.
A quadratic trend $X_t = a + b\,t + c\,t^2$ is reduced to a constant
by $\Delta^2$:
$\Delta^2 X_t = 2c$, which is stationary.
\end{exbox}

\begin{defbox}[Integration order $I(d)$]
$X_t \sim I(d)$ if $\Delta^d X_t$ is stationary but
$\Delta^{d - 1} X_t$ is not. White noise is $I(0)$; a random walk is
$I(1)$; a random walk with a stochastic trend is often $I(2)$.
\end{defbox}

% =========================================================
\section{Unit-root tests}
\label{sec:unit-roots}

\subsection{Augmented Dickey--Fuller (ADF) test}

\begin{defbox}[ADF regression]
\label{def:adf}
Starting from an AR(1) $X_t = \rho X_{t-1} + \varepsilon_t$,
substitute $\rho = 1 + \gamma$ and subtract $X_{t-1}$:
\[
  \Delta X_t = \gamma\, X_{t-1} + \varepsilon_t,
\]
augmented with $p$ lags of $\Delta X_t$ to whiten the error:
\[
  \Delta X_t = \gamma\, X_{t-1}
             + \sum_{i = 1}^{p} \alpha_i\, \Delta X_{t-i}
             + \varepsilon_t.
\]
\end{defbox}

\begin{thbox}[ADF test]
\label{th:adf}
\begin{itemize}
  \item $H_0 : \gamma = 0$ — unit root, $X_t$ is non-stationary.
  \item $H_1 : \gamma < 0$ — $X_t$ is stationary.
\end{itemize}
The $t$-statistic of $\widehat\gamma$ \emph{does not} follow a
Student's $t$ under $H_0$; its asymptotic distribution is the
\emph{Dickey--Fuller distribution}, with 5\% critical value
$\approx -2.86$ for the constant-only specification (and
$\approx -3.41$ with constant and trend).
\end{thbox}

\begin{rembox}[Three flavours]
\begin{itemize}
  \item \emph{No constant, no trend} (rare): $X_t$ should fluctuate
        around zero.
  \item \emph{Constant only} (most common): $X_t$ should fluctuate
        around a non-zero level.
  \item \emph{Constant and trend}: $X_t$ has a deterministic linear
        trend.
\end{itemize}
The choice affects the critical values; the Python function
\texttt{adfuller} accepts a \texttt{regression} argument.
\end{rembox}

\subsection{KPSS test}

\begin{defbox}[KPSS test, Kwiatkowski et al.\ 1992]
\label{def:kpss}
\[
  H_0 : \text{stationarity (around a level or a trend)},
  \qquad
  H_1 : \text{unit-root non-stationarity}.
\]
The test statistic is built from the partial sums of OLS residuals
versus a long-run variance estimator. 5\% critical values: $0.463$
for level stationarity, $0.146$ for trend stationarity. Reject $H_0$
if the statistic \emph{exceeds} the critical value.
\end{defbox}

\begin{rembox}[Combining ADF and KPSS]
The two tests have \emph{opposite} null hypotheses, which clarifies
ambiguous cases:
\begin{center}
\begin{tabular}{l|cc}
                              & ADF rejects        & ADF does not reject \\
\hline
KPSS rejects                  & contradiction      & non-stationary \\
KPSS does not reject          & stationary         & inconclusive (small $T$) \\
\end{tabular}
\end{center}
The contradiction case typically arises when there is a structural
break.
\end{rembox}

\subsection{Phillips--Perron test (mention)}

\begin{rembox}
The \emph{Phillips--Perron} (PP) test is an alternative to ADF that
adjusts the test statistic for heteroscedasticity and autocorrelation
non-parametrically (instead of via lagged differences). Same null
hypothesis as ADF, similar critical values.
\end{rembox}

% =========================================================
\section{ARIMA(p, d, q)}
\label{sec:arima}

\begin{defbox}[ARIMA model]
\label{def:arima}
$X_t \sim \ARIMA{p}{d}{q}$ if its $d$-th difference is a stationary
$\ARMA{p}{q}$:
\[
  \phi(B)\, (1 - B)^d\, X_t \;=\; \theta(B)\, \varepsilon_t,
  \qquad \varepsilon_t \sim \WN(0, \sigma^2).
\]
\end{defbox}

\begin{rembox}[Practical recipe]
\begin{enumerate}
  \item Plot $X_t$; identify visual non-stationarity.
  \item Test (ADF + KPSS); difference until stationary; choose $d$.
  \item ACF/PACF of $\Delta^d X_t$ $\Rightarrow$ candidate $(p, q)$.
  \item MLE estimate of the ARMA(p, q) coefficients.
  \item Ljung--Box on residuals; if not white, increase $p$ or $q$.
  \item Forecast, integrate back.
\end{enumerate}
\end{rembox}

\begin{exbox}[ARIMA(1, 1, 1)]
\label{ex:arima-111}
$\phi(B) = 1 - \phi B$, $\theta(B) = 1 + \theta B$, $d = 1$:
\[
  (1 - \phi B)(1 - B) X_t = (1 + \theta B)\, \varepsilon_t.
\]
With $W_t = \Delta X_t$:
$W_t = \phi W_{t-1} + \varepsilon_t + \theta \varepsilon_{t-1}$,
i.e.\ $W_t \sim \ARMA{1}{1}$. Starting from a known initial value
$X_0$, integrate back: $X_t = X_0 + \sum_{s = 1}^{t} W_s$.
\end{exbox}

% =========================================================
\section{Seasonal ARIMA (SARIMA)}
\label{sec:sarima}

When the series exhibits seasonality of period $s$, classical
differencing $(1 - B)^d$ alone is not enough — we also apply
seasonal differencing $(1 - B^s)^D$ and seasonal ARMA polynomials.

\begin{defbox}[SARIMA(p, d, q)(P, D, Q)$_s$]
\label{def:sarima}
\[
  \phi(B)\, \Phi(B^s)\,
  (1 - B)^d\, (1 - B^s)^D\, X_t
  \;=\;
  \theta(B)\, \Theta(B^s)\, \varepsilon_t,
\]
where $\Phi(z) = 1 - \Phi_1 z - \dots - \Phi_P z^P$ and
$\Theta(z) = 1 + \Theta_1 z + \dots + \Theta_Q z^Q$ are the
\emph{seasonal} AR and MA polynomials.
\end{defbox}

\begin{exbox}[SARIMA(1, 1, 1)(1, 1, 1)$_{12}$]
\[
  (1 - \phi B)(1 - \Phi B^{12})(1 - B)(1 - B^{12}) X_t
  = (1 + \theta B)(1 + \Theta B^{12}) \varepsilon_t.
\]
This model is the workhorse for monthly time series with both
short-term persistence and yearly seasonality (sales, climate, energy
demand).
\end{exbox}

\begin{rembox}[Identification]
After the differencing step, look at the ACF/PACF of
$(1 - B)^d (1 - B^s)^D X_t$:
\begin{itemize}
  \item significant spikes at lags $s, 2s, 3s, \dots$ in the ACF
        $\Rightarrow$ seasonal MA component (order $Q$);
  \item spikes at $s, 2s, \dots$ in the PACF $\Rightarrow$ seasonal
        AR component (order $P$);
  \item spikes at small lags $\Rightarrow$ ordinary $p, q$.
\end{itemize}
\end{rembox}

% =========================================================
\section*{Summary}
\addcontentsline{toc}{section}{Summary}

\noindent\textbf{Key takeaways.}
\begin{itemize}
  \item Most real series are non-stationary; \emph{spurious
        regression} can fool naive applications of OLS or ARMA.
  \item Distinguish \emph{trend stationary} (detrend) from
        \emph{difference stationary} (difference) using the ADF and
        KPSS tests.
  \item The backshift $B$ and difference $\Delta$ operators give a
        compact algebra for non-stationary models; integration order
        $I(d)$ counts the number of differencings needed.
  \item \emph{ARIMA(p, d, q)} = ARMA on the $d$-th difference;
        \emph{SARIMA} adds seasonal AR, MA and differencing.
  \item In Python, \texttt{statsmodels.tsa.stattools.adfuller} and
        \texttt{kpss}, plus \texttt{SARIMAX}, cover the whole
        workflow.
\end{itemize}


% =========================================================
% Chapter 5 -- ARMA models
% Author : Aymen Ben Brik (aymen.benbrik@esprit.tn)
% =========================================================

\chapter{ARMA models}
\label{ch:arma}

\section*{Introduction}
\addcontentsline{toc}{section}{Introduction}

Chapter~\ref{ch:stationarity} introduced stationarity and the
sample ACF/PACF; Chapter~\ref{ch:nonstationarity} showed how to
turn a non-stationary series into a stationary one by differencing.
This chapter assembles the central modelling family for stationary
time series: the \emph{autoregressive moving-average} processes,
ARMA$(p, q)$. We give their precise definitions, state the
stationarity and invertibility conditions in terms of polynomial
roots, derive their ACF/PACF \emph{signatures}, and detail the
practical workflow --- the Box--Jenkins method --- for identification,
estimation, diagnostics and forecasting.

% =========================================================
\section{Autoregressive process AR$(p)$}
\label{sec:ar}

\begin{defbox}[Autoregressive process AR$(p)$]
\label{def:ar}
A process $(X_t)_{t \in \Z}$ is autoregressive of order $p$ if
\[
  X_t = \phi_1 X_{t-1} + \phi_2 X_{t-2} + \dots + \phi_p X_{t-p}
        + \varepsilon_t,
  \qquad \varepsilon_t \sim \WN(0, \sigma^2),
\]
with $\phi_p \neq 0$. Equivalently, with $\Phi(z) = 1 - \phi_1 z -
\dots - \phi_p z^p$,
\[
  \Phi(B)\, X_t = \varepsilon_t.
\]
\end{defbox}

\begin{propbox}[Stationarity of AR$(p)$]
\label{prop:ar-stationary}
Equation $\Phi(B)\, X_t = \varepsilon_t$ admits a unique causal
stationary solution if and only if every root $z_i$ of
$\Phi(z) = 0$ satisfies $\abs{z_i} > 1$. In that case
\[
  X_t = \Phi^{-1}(B)\, \varepsilon_t
      = \sum_{j = 0}^{\infty} \psi_j\, \varepsilon_{t-j},
\]
with $\sum_{j} \abs{\psi_j} < \infty$.
\end{propbox}

\begin{exbox}[AR$(1)$ moments]
\label{ex:ar1}
For $X_t = \phi X_{t-1} + \varepsilon_t$ with $\abs\phi < 1$,
iterating gives
$X_t = \sum_{j \geq 0} \phi^j\, \varepsilon_{t - j}$, hence
\[
  \E[X_t] = 0,
  \qquad
  \Var(X_t) = \frac{\sigma^2}{1 - \phi^2},
  \qquad
  \rho(h) = \phi^{\abs h}.
\]
The ACF \textbf{decays geometrically} and the PACF \textbf{cuts off}
after lag~$1$ ($\alpha(1) = \phi$, $\alpha(h) = 0$ for $h \geq 2$).
\end{exbox}

\begin{propbox}[Yule--Walker equations]
\label{prop:yw}
Let $X_t \sim \AR{p}$ be stationary. Multiplying the equation by
$X_{t - h}$ and taking expectations gives, for $h = 1, \dots, p$,
\[
  \rho(h) = \phi_1\, \rho(h - 1) + \phi_2\, \rho(h - 2)
          + \dots + \phi_p\, \rho(h - p).
\]
In matrix form $\bm R \bm \phi = \bm r$, the system is invertible
under the stationarity condition; replacing $\rho$ by the sample
ACF $\widehat\rho$ gives the Yule--Walker estimator
$\widehat\phi_{\text{YW}}$.
\end{propbox}

\begin{exbox}[AR$(2)$ stationarity triangle]
For $X_t = \phi_1 X_{t-1} + \phi_2 X_{t-2} + \varepsilon_t$, the
roots of $1 - \phi_1 z - \phi_2 z^2 = 0$ lie outside the unit disk
iff
\[
  \phi_1 + \phi_2 < 1,
  \qquad
  \phi_2 - \phi_1 < 1,
  \qquad
  -1 < \phi_2 < 1.
\]
This is the Schur stability condition for the AR(2) polynomial.
\end{exbox}

% =========================================================
\section{Moving-average process MA$(q)$}
\label{sec:ma}

\begin{defbox}[Moving-average process MA$(q)$]
\label{def:ma}
A process $(X_t)_{t \in \Z}$ is moving-average of order $q$ if
\[
  X_t = \varepsilon_t + \theta_1 \varepsilon_{t-1}
        + \dots + \theta_q \varepsilon_{t-q},
  \qquad \varepsilon_t \sim \WN(0, \sigma^2),
\]
with $\theta_q \neq 0$. Equivalently, $X_t = \Theta(B)\,
\varepsilon_t$ with $\Theta(z) = 1 + \theta_1 z + \dots + \theta_q
z^q$.
\end{defbox}

\begin{propbox}[Moments and ACF of MA$(q)$]
\label{prop:ma-acf}
Every MA$(q)$ is (weakly) stationary, with
\[
  \E[X_t] = 0,
  \qquad
  \Var(X_t) = \sigma^2 (1 + \theta_1^2 + \dots + \theta_q^2),
\]
and autocovariance
\[
  \gamma(h) =
  \begin{cases}
    \sigma^2 \bigl(\theta_h + \theta_1 \theta_{h+1} + \dots
                 + \theta_{q-h} \theta_q\bigr)
                                & 1 \leq h \leq q, \\[4pt]
    0                           & h > q,
  \end{cases}
\]
with $\theta_0 = 1$. In particular $\rho(h) = 0$ for $h > q$:
the ACF \textbf{cuts off} at lag~$q$.
\end{propbox}

\begin{defbox}[Invertibility]
\label{def:invert}
An MA$(q)$ is \emph{invertible} if it admits an AR$(\infty)$
representation
\[
  \varepsilon_t = X_t - \pi_1 X_{t-1} - \pi_2 X_{t-2} - \dots,
  \qquad \sum_{j} \abs{\pi_j} < \infty.
\]
This holds if and only if every root of $\Theta(z) = 0$ lies
\emph{outside} the closed unit disk.
\end{defbox}

\begin{rembox}[Why invertibility matters]
Without invertibility, two parameter sets give the same ACF and so
are indistinguishable from the data. The MA$(1)$ models
$X_t = \varepsilon_t + \theta \varepsilon_{t-1}$ and
$X_t = \tilde\varepsilon_t + (1/\theta)\, \tilde\varepsilon_{t-1}$
have identical autocovariance functions; we conventionally choose
the invertible one, $\abs\theta < 1$.
\end{rembox}

% =========================================================
\section{ARMA$(p, q)$}
\label{sec:arma}

\begin{defbox}[ARMA$(p, q)$]
\label{def:arma}
A process $(X_t)$ is ARMA$(p, q)$ if it satisfies
\[
  \Phi(B)\, X_t = \Theta(B)\, \varepsilon_t,
  \qquad \varepsilon_t \sim \WN(0, \sigma^2),
\]
with $\Phi(z) = 1 - \phi_1 z - \dots - \phi_p z^p$,
$\Theta(z) = 1 + \theta_1 z + \dots + \theta_q z^q$, and $\Phi$,
$\Theta$ sharing no common root.
\end{defbox}

\begin{propbox}[Stationarity, invertibility and existence]
\label{prop:arma}
The equation $\Phi(B) X_t = \Theta(B) \varepsilon_t$ admits a
unique causal stationary invertible solution if and only if
\[
  \abs{z_i} > 1 \text{ for every root of } \Phi,
  \qquad
  \abs{z_j} > 1 \text{ for every root of } \Theta.
\]
The solution is then
\[
  X_t = \Phi^{-1}(B)\, \Theta(B)\, \varepsilon_t
      = \sum_{j = 0}^{\infty} \psi_j\, \varepsilon_{t-j}.
\]
\end{propbox}

\begin{rembox}
The MA$(\infty)$ coefficients $(\psi_j)$ are obtained by polynomial
division of $\Theta(z)$ by $\Phi(z)$ around $z = 0$. They
control both the autocovariance and the forecast variance
(Section~\ref{sec:forecasting}).
\end{rembox}

% =========================================================
\section{Box--Jenkins identification}
\label{sec:identification}

The choice of $(p, q)$ is guided by the sample ACF and PACF, whose
\emph{signatures} differ across model classes:

\begin{center}
\renewcommand{\arraystretch}{1.4}
\begin{tabular}{lll}
\toprule
Model        & ACF $\rho(h)$              & PACF $\alpha(h)$ \\
\midrule
AR$(p)$      & exponential decay          & cuts off after lag $p$ \\
MA$(q)$      & cuts off after lag $q$     & exponential decay \\
ARMA$(p, q)$ & exponential decay (after $q$) & exponential decay (after $p$) \\
\bottomrule
\end{tabular}
\end{center}

\begin{rembox}[Bartlett bands]
``Cuts off'' is read as $\widehat\rho(h)$ falling \emph{inside} the
$\pm 1.96 / \sqrt T$ band for $h$ larger than the cutoff lag.
``Decays'' means the values stay outside the band but shrink in
absolute value.
\end{rembox}

\begin{exbox}[Reading a sample ACF/PACF]
On a series of length $T = 200$, the sample PACF shows
$\widehat\alpha(1) = 0.62$, $\widehat\alpha(2) = -0.04$,
$\widehat\alpha(3) = 0.02$ and the sample ACF decays geometrically.
The Bartlett band is $\pm 1.96 / \sqrt{200} \approx \pm 0.139$, so
$\widehat\alpha(h)$ is inside the band for $h \geq 2$:
the PACF cuts off at lag~$1$, suggesting AR$(1)$.
\end{exbox}

% =========================================================
\section{Estimation}
\label{sec:arma-estimation}

\subsection{Yule--Walker (AR only)}

For pure AR$(p)$, plug $\widehat\rho$ into
Proposition~\ref{prop:yw} to obtain the Yule--Walker estimate.
This is method-of-moments: \emph{consistent} but not asymptotically
efficient.

\subsection{Conditional sum of squares (CSS)}

Conditioning on the first $p$ observations and assuming
$\varepsilon_t = 0$ for $t \leq 0$, the CSS estimator is
\[
  \widehat{\bm\theta}_{\text{CSS}}
  = \arg\min_{\bm\theta} \sum_{t = p + 1}^{T}
        \widehat\varepsilon_t(\bm\theta)^2,
\]
where $\widehat\varepsilon_t$ is computed recursively from the
ARMA equation. CSS is fast and equivalent to MLE in large samples
under Gaussian innovations.

\subsection{Maximum likelihood (MLE)}

Under $\varepsilon_t \iid \mathcal N(0, \sigma^2)$, the joint
distribution of $(X_1, \dots, X_T)$ is multivariate normal with
covariance matrix $\Gamma_T(\bm\theta)$. The log-likelihood
\[
  \log \mathcal L(\bm\theta, \sigma^2)
  = -\tfrac{T}{2}\log(2\pi)
    -\tfrac{1}{2}\log \det \Gamma_T(\bm\theta)
    -\tfrac{1}{2} \bm X^\top \Gamma_T(\bm\theta)^{-1} \bm X
\]
is maximised numerically. The standard implementation in
\texttt{statsmodels} uses a state-space representation and the
Kalman filter, which scale linearly in $T$ and provide asymptotic
standard errors.

\begin{thbox}[Asymptotic distribution of the MLE]
Under regularity assumptions (stationarity, invertibility, no
common roots), as $T \to \infty$,
\[
  \sqrt{T}\, (\widehat{\bm\theta}_{\text{MLE}} - \bm\theta_0)
  \;\xrightarrow{d}\; \mathcal N(0, I(\bm\theta_0)^{-1}),
\]
where $I(\bm\theta_0)$ is the Fisher information matrix evaluated
at the true parameter $\bm\theta_0$.
\end{thbox}

% =========================================================
\section{Model selection}
\label{sec:selection}

\begin{defbox}[Information criteria]
\label{def:aic-bic}
For a fitted ARMA$(p, q)$ with $k = p + q + 1$ free parameters
(including $\sigma^2$),
\[
  \mathrm{AIC} = -2 \log \widehat{\mathcal L} + 2 k,
  \qquad
  \mathrm{BIC} = -2 \log \widehat{\mathcal L} + k \log T.
\]
The chosen order minimises the criterion.
\end{defbox}

\begin{rembox}[AIC vs BIC]
AIC favours predictive accuracy and tends to over-fit on small
samples; BIC penalises complexity more and is consistent for
the true order under standard assumptions. When the two disagree,
prefer the BIC choice for parsimony and confirm with the residual
diagnostics.
\end{rembox}

% =========================================================
\section{Residual diagnostics}
\label{sec:diagnostics}

\begin{rembox}[Adequacy checks]
A well-specified ARMA model produces residuals
$\widehat\varepsilon_t = X_t - \widehat X_{t \mid t-1}$ that are
\emph{indistinguishable from white noise}. Standard checks:
\begin{itemize}
  \item time plot of $\widehat\varepsilon_t$ -- mean stable around
        zero, variance constant;
  \item sample ACF of $\widehat\varepsilon_t$ -- spikes inside
        Bartlett bands;
  \item Ljung--Box test;
  \item QQ-plot vs $\mathcal N(0, 1)$ for normality (Jarque--Bera
        test).
\end{itemize}
\end{rembox}

\begin{thbox}[Ljung--Box test on residuals]
\label{th:lb-residuals}
\[
  Q_{\text{LB}}(m)
  = T(T + 2) \sum_{h = 1}^{m}
        \frac{\widehat\rho_{\widehat\varepsilon}(h)^2}{T - h}.
\]
Under the null that $\widehat\varepsilon_t$ is white noise, after
fitting an ARMA$(p, q)$,
\[
  Q_{\text{LB}}(m) \;\xrightarrow{d}\; \chi^2_{m - p - q}.
\]
Reject the null if $Q_{\text{LB}}(m)$ exceeds the
$\chi^2_{m - p - q, 1 - \alpha}$ quantile.
\end{thbox}

% =========================================================
\section{Forecasting}
\label{sec:forecasting}

\begin{defbox}[Best linear forecast]
\label{def:forecast}
The minimum mean-squared-error linear forecast at horizon $h$,
given the past $\mathcal F_T = \sigma(X_T, X_{T-1}, \dots)$, is
\[
  \widehat X_{T + h \mid T}
  = \E\bigl[X_{T + h} \mid \mathcal F_T \bigr].
\]
The forecast error is
$e_{T + h} = X_{T + h} - \widehat X_{T + h \mid T}$.
\end{defbox}

\begin{propbox}[Forecast error variance]
\label{prop:forecast-var}
With the MA$(\infty)$ representation
$X_t = \sum_{j \geq 0} \psi_j\, \varepsilon_{t - j}$,
\[
  e_{T + h} = \sum_{j = 0}^{h - 1} \psi_j\, \varepsilon_{T + h - j},
\]
hence
\[
  \E\bigl[e_{T + h}\bigr] = 0,
  \qquad
  \Var\bigl(e_{T + h}\bigr)
  = \sigma^2 \sum_{j = 0}^{h - 1} \psi_j^2.
\]
\end{propbox}

\begin{exbox}[Forecast for AR$(1)$]
With $X_t = \phi X_{t-1} + \varepsilon_t$, $\abs\phi < 1$:
\[
  \widehat X_{T + h \mid T} = \phi^h X_T,
  \qquad
  \Var(e_{T + h}) = \sigma^2 \frac{1 - \phi^{2h}}{1 - \phi^2}.
\]
As $h \to \infty$, $\widehat X_{T + h \mid T} \to 0$ (the
unconditional mean) and the variance tends to
$\sigma^2 / (1 - \phi^2)$ (the unconditional variance).
\end{exbox}

\begin{rembox}[Confidence intervals]
Under Gaussian innovations, the $1 - \alpha$ confidence interval is
\[
  \widehat X_{T + h \mid T}
  \;\pm\; z_{1 - \alpha / 2}\, \sigma\,
        \sqrt{\sum_{j = 0}^{h - 1} \psi_j^2}.
\]
For stationary ARMA, the width converges to a finite limit; for
ARIMA (Chapter~\ref{ch:nonstationarity}), it grows without bound.
\end{rembox}

% =========================================================
\section*{Summary}
\addcontentsline{toc}{section}{Summary}

\noindent\textbf{Key takeaways.}
\begin{itemize}
  \item AR$(p)$ depends on past values; MA$(q)$ depends on past
        shocks; ARMA$(p, q)$ combines both.
  \item Stationarity = roots of $\Phi$ outside the unit disk;
        invertibility = roots of $\Theta$ outside the unit disk.
  \item ACF/PACF signatures: AR cuts the PACF, MA cuts the ACF,
        ARMA shows double exponential decay.
  \item Estimation: Yule--Walker for pure AR; CSS or MLE for
        general ARMA. \texttt{statsmodels.tsa.arima.ARIMA} provides
        both.
  \item Model selection by AIC / BIC; validate with Ljung--Box on
        residuals.
  \item Forecasts widen with horizon according to
        $\Var(e_{T + h}) = \sigma^2 \sum_{j = 0}^{h - 1} \psi_j^2$.
\end{itemize}


% =========================================================
% Chapter 6 -- ARCH and GARCH models
% Author : Aymen Ben Brik (aymen.benbrik@esprit.tn)
% =========================================================

\chapter{ARCH and GARCH models}
\label{ch:arch-garch}

\section*{Introduction}
\addcontentsline{toc}{section}{Introduction}

The ARMA family of Chapter~\ref{ch:arma} captures the dynamics of
the \emph{conditional mean}, but assumes a constant variance. On
financial returns, on macroeconomic surprises and on many physical
measurements, this assumption is violated: large absolute deviations
cluster in time, and the variance itself is predictable.
This chapter introduces the \emph{autoregressive conditional
heteroskedasticity} (ARCH) family that models the conditional
variance directly, generalises it to GARCH$(p, q)$, derives the
stationarity conditions, presents the maximum-likelihood estimator,
and applies the model to volatility forecasting and Value-at-Risk.

% =========================================================
\section{Stylised facts of financial returns}
\label{sec:stylised}

\begin{rembox}[Empirical regularities]
For daily log-returns $r_t = \log(P_t / P_{t-1})$ on broad equity
indices, currencies and commodities:
\begin{enumerate}
  \item $\E[r_t] \approx 0$ (small drift).
  \item ACF of $r_t$: nearly flat -- returns are approximately
        serially uncorrelated.
  \item ACF of $r_t^2$ (or $|r_t|$): \emph{strongly} positive at
        many lags -- volatility is autocorrelated.
  \item Marginal kurtosis $\widehat\kappa > 3$: heavier tails than
        Gaussian.
  \item Leverage effect: negative shocks raise future volatility
        more than positive shocks of the same size.
\end{enumerate}
\end{rembox}

\begin{rembox}[Conditional vs.\ unconditional]
Let $\mathcal F_{t-1} = \sigma(r_{t-1}, r_{t-2}, \dots)$ be the
information set up to $t - 1$.
\begin{itemize}
  \item \emph{Unconditional variance}:
        $\sigma^2 = \Var(r_t)$, the long-run, average risk.
  \item \emph{Conditional variance}:
        $\sigma_t^2 = \Var(r_t \mid \mathcal F_{t-1})$, the
        short-run risk knowing the past.
\end{itemize}
ARCH and GARCH models \emph{specify $\sigma_t^2$} while keeping
$\sigma^2$ finite -- the process is unconditionally homoskedastic
but conditionally heteroskedastic.
\end{rembox}

\begin{defbox}[Mean--variance decomposition]
\label{def:mean-var}
A second-order martingale-difference innovation
$\varepsilon_t = r_t - \mu_t$ is written
\[
  \varepsilon_t = \sigma_t\, z_t,
  \qquad
  z_t \iid \mathcal D(0, 1),
  \qquad
  \E[\varepsilon_t \mid \mathcal F_{t-1}] = 0,
  \quad
  \Var(\varepsilon_t \mid \mathcal F_{t-1}) = \sigma_t^2.
\]
$\mathcal D$ is the \emph{innovation distribution}, typically
$\mathcal N(0, 1)$ or Student-$t$ with low degrees of freedom.
\end{defbox}

% =========================================================
\section{ARCH$(q)$ model}
\label{sec:arch}

\begin{defbox}[ARCH$(q)$, Engle 1982]
\label{def:arch}
The \emph{autoregressive conditional heteroskedasticity} model of
order $q$ specifies
\[
  \varepsilon_t = \sigma_t z_t,
  \qquad
  \sigma_t^2 = \alpha_0
              + \sum_{i = 1}^{q} \alpha_i\, \varepsilon_{t-i}^2,
\]
with $z_t \iid \mathcal N(0, 1)$, $\alpha_0 > 0$, and
$\alpha_i \geq 0$ for $i = 1, \dots, q$.
\end{defbox}

\begin{propbox}[Properties of ARCH$(q)$]
\label{prop:arch}
Assume the covariance-stationarity condition
$\sum_{i = 1}^{q} \alpha_i < 1$. Then:
\begin{enumerate}
  \item $\E[\varepsilon_t] = 0$ and $(\varepsilon_t)$ is a white
        noise (uncorrelated in time, but not independent).
  \item $\E[\varepsilon_t^2] = \alpha_0 / (1 - \sum \alpha_i)
        =: \sigma^2$.
  \item $(\varepsilon_t^2)$ admits an AR$(q)$ representation
        $\varepsilon_t^2 = \alpha_0
                          + \sum \alpha_i \varepsilon_{t-i}^2 + \nu_t$,
        with $\nu_t = \varepsilon_t^2 - \sigma_t^2$ a martingale
        difference.
  \item Excess kurtosis: under Gaussian innovations
        ARCH$(1)$ has kurtosis
        $3\,(1 - \alpha_1^2) / (1 - 3 \alpha_1^2) > 3$ when
        $\alpha_1^2 < 1/3$.
\end{enumerate}
\end{propbox}

\begin{rembox}[Why squares?]
Squaring past residuals makes the response \emph{sign-blind}: a
$-3 \sigma$ surprise and a $+3 \sigma$ surprise both raise
$\sigma_{t+1}^2$. To allow asymmetric responses (leverage), one
moves to GJR-GARCH or EGARCH (Section~\ref{sec:extensions}).
\end{rembox}

% =========================================================
\section{GARCH$(p, q)$ model}
\label{sec:garch}

In practice, ARCH$(q)$ requires many lags $q$ to capture the
slow decay of squared-return autocorrelations, leading to long,
poorly identified parameter vectors. Bollerslev (1986) generalised
the model by including past conditional variances themselves.

\begin{defbox}[GARCH$(p, q)$, Bollerslev 1986]
\label{def:garch}
\[
  \sigma_t^2 = \alpha_0
              + \sum_{i = 1}^{q} \alpha_i\, \varepsilon_{t-i}^2
              + \sum_{j = 1}^{p} \beta_j\, \sigma_{t-j}^2,
\]
with $\alpha_0 > 0$, $\alpha_i, \beta_j \geq 0$.
\end{defbox}

\begin{propbox}[Stationarity of GARCH$(p, q)$]
\label{prop:garch-stat}
$(\varepsilon_t)$ admits a covariance-stationary solution with
finite unconditional variance if and only if
\[
  \sum_{i = 1}^{q} \alpha_i + \sum_{j = 1}^{p} \beta_j \;<\; 1.
\]
The unconditional variance is then
\[
  \sigma^2 = \E[\varepsilon_t^2]
          = \frac{\alpha_0}{1 - \sum_i \alpha_i - \sum_j \beta_j}.
\]
\end{propbox}

\begin{exbox}[GARCH$(1, 1)$: the workhorse]
\label{ex:garch11}
\[
  \sigma_t^2 = \alpha_0
             + \alpha_1\, \varepsilon_{t-1}^2
             + \beta_1\, \sigma_{t-1}^2.
\]
On daily equity returns, typical estimates are $\alpha_1 \approx
0.05$--$0.10$, $\beta_1 \approx 0.85$--$0.92$,
$\alpha_1 + \beta_1 \approx 0.95$--$0.99$. Volatility is
\emph{very persistent} but still mean-reverting.
\end{exbox}

\begin{rembox}[IGARCH and persistence]
The boundary case $\alpha_1 + \beta_1 = 1$ defines an
\emph{IGARCH}: shocks have a permanent effect on the conditional
variance. The unconditional variance is then infinite (the model
is strictly stationary but not covariance-stationary). Estimated
$\alpha_1 + \beta_1$ very close to $1$ is a sign that the assumed
GARCH$(1, 1)$ may be missing a structural break.
\end{rembox}

% =========================================================
\section{Maximum-likelihood estimation}
\label{sec:garch-mle}

\begin{defbox}[Conditional Gaussian likelihood]
\label{def:garch-mle}
Assuming $z_t \iid \mathcal N(0, 1)$ and $r_t = \mu + \varepsilon_t$,
the conditional log-likelihood is
\[
  \log \mathcal L(\bm\theta)
  = -\tfrac{T}{2}\log(2 \pi)
    -\tfrac{1}{2} \sum_{t = 1}^{T}
        \Bigl(
            \log \sigma_t^2(\bm\theta)
            + \frac{(r_t - \mu)^2}{\sigma_t^2(\bm\theta)}
        \Bigr),
\]
with $\sigma_t^2$ computed recursively from the GARCH equation.
$\bm\theta = (\mu, \alpha_0, \alpha_1, \dots, \alpha_q,
\beta_1, \dots, \beta_p)$.
\end{defbox}

\begin{thbox}[Asymptotic properties of the QMLE]
Under regularity conditions (positivity, stationarity, bounded
fourth moment), the Gaussian quasi-MLE
$\widehat{\bm\theta}_{\text{QMLE}}$ is consistent and
asymptotically normal,
\[
  \sqrt{T}\,(\widehat{\bm\theta} - \bm\theta_0)
  \;\xrightarrow{d}\; \mathcal N(0, \Sigma),
\]
\emph{even if} the true innovation distribution is non-Gaussian
(e.g.\ Student-$t$). The asymptotic variance $\Sigma$ takes a
sandwich form $\Sigma = J^{-1} I J^{-1}$.
\end{thbox}

\begin{rembox}[Practical numerics]
\begin{itemize}
  \item Returns are typically rescaled (multiplied by $100$ to work
        in percentage units) for numerical stability.
  \item Initial $\sigma_1^2$ is set to the sample variance, or to
        the unconditional value $\alpha_0 / (1 - \alpha_1 -
        \beta_1)$ at the current parameters.
  \item In Python, the \texttt{arch} package
        (\texttt{arch\_model}) does the rescaling, parameter
        constraints and likelihood maximisation automatically.
\end{itemize}
\end{rembox}

% =========================================================
\section{Diagnostics}
\label{sec:garch-diag}

\subsection{Engle's LM test for ARCH effects}

\begin{defbox}[Engle's LM test]
\label{def:engle-lm}
Apply on the residual $\widehat\varepsilon_t$ of a mean model
(constant or ARMA). Run the auxiliary regression
\[
  \widehat\varepsilon_t^2
  = \gamma_0 + \gamma_1 \widehat\varepsilon_{t-1}^2
            + \dots + \gamma_q \widehat\varepsilon_{t-q}^2 + u_t,
\]
get its $R^2$. Test
\[
  H_0: \gamma_1 = \dots = \gamma_q = 0,
  \qquad
  LM = T \cdot R^2 \;\xrightarrow{d}\; \chi^2_q.
\]
Reject $H_0$ at level $\alpha$ if $LM > \chi^2_{q, 1 - \alpha}$.
\end{defbox}

\begin{rembox}[How to use it]
Run Engle's LM \emph{before} fitting GARCH (to confirm there is an
ARCH effect to model) and \emph{after} (on the standardised
residuals $\widehat z_t = \widehat\varepsilon_t / \widehat\sigma_t$,
to confirm the model has captured it).
\end{rembox}

\subsection{Standardised residual diagnostics}

\begin{rembox}
A well-fitted GARCH produces standardised residuals that look like
white noise both in level and in square:
\begin{itemize}
  \item ACF of $\widehat z_t$ inside Bartlett bands;
  \item ACF of $\widehat z_t^2$ inside Bartlett bands;
  \item Engle LM on $\widehat z_t^2$ does not reject;
  \item marginal distribution close to $\mathcal D(0, 1)$ (QQ-plot,
        Kolmogorov--Smirnov).
\end{itemize}
Persistent failures in $\widehat z_t^2$ suggest a higher-order GARCH
or a Student-$t$ innovation; failures only in $\widehat z_t^-$
suggest a leverage extension (GJR, EGARCH).
\end{rembox}

% =========================================================
\section{Volatility forecasting and Value-at-Risk}
\label{sec:garch-forecast}

\begin{propbox}[GARCH$(1, 1)$ multi-step forecast]
\label{prop:garch-forecast}
Let $\widehat\sigma_T^2$ and $\widehat\varepsilon_T$ be the values
obtained at the end of the sample. Then
\[
  \widehat\sigma_{T+1}^2
  = \alpha_0 + \alpha_1 \widehat\varepsilon_T^2
             + \beta_1 \widehat\sigma_T^2,
\]
and for $h \geq 2$,
\[
  \widehat\sigma_{T+h}^2
  = \sigma^2
  + (\alpha_1 + \beta_1)^{h-1}
        \bigl(\widehat\sigma_{T+1}^2 - \sigma^2\bigr),
\]
where $\sigma^2 = \alpha_0 / (1 - \alpha_1 - \beta_1)$ is the
unconditional variance. The forecast \emph{mean-reverts} to
$\sigma^2$ at rate $\alpha_1 + \beta_1$.
\end{propbox}

\begin{defbox}[Conditional Value-at-Risk]
\label{def:var}
At horizon $1$ and confidence $1 - \alpha$ (typically $\alpha = 1\%$
or $5\%$), assuming Gaussian innovations,
\[
  \mathrm{VaR}_{T+1}(\alpha)
  = -\bigl(\widehat\mu + z_\alpha\, \widehat\sigma_{T+1}\bigr),
\]
with $z_\alpha = \Phi^{-1}(\alpha)$ ($z_{0.01} = -2.326$,
$z_{0.05} = -1.645$). The interpretation: with probability
$1 - \alpha$, the loss does not exceed $\mathrm{VaR}_{T+1}(\alpha)$.
\end{defbox}

\begin{rembox}[Static vs dynamic VaR]
The classical static VaR uses the \emph{sample} standard deviation
$\widehat\sigma$. The GARCH-VaR uses
$\widehat\sigma_{T+1}$. In calm regimes, $\widehat\sigma_{T+1} <
\widehat\sigma$ and the dynamic VaR is tighter; in turbulent
regimes (just after a market shock), $\widehat\sigma_{T+1} >
\widehat\sigma$ and the dynamic VaR is wider, accurately reflecting
the elevated short-term risk.
\end{rembox}

% =========================================================
\section{Extensions (overview)}
\label{sec:extensions}

\begin{rembox}[EGARCH (Nelson 1991)]
The exponential GARCH models $\log \sigma_t^2$, removing the
positivity constraints on the parameters and naturally allowing
asymmetric responses:
\[
  \log \sigma_t^2
  = \omega + \alpha\, \bigl(|z_{t-1}| - \E |z_{t-1}|\bigr)
          + \gamma\, z_{t-1}
          + \beta\, \log \sigma_{t-1}^2.
\]
$\gamma < 0$ is the typical sign of the leverage effect on equity
returns.
\end{rembox}

\begin{rembox}[GJR-GARCH (Glosten--Jagannathan--Runkle 1993)]
Adds a threshold term that activates only on negative shocks:
\[
  \sigma_t^2
  = \alpha_0
  + (\alpha_1 + \gamma\, \mathbb 1_{\varepsilon_{t-1} < 0})
       \varepsilon_{t-1}^2
  + \beta_1\, \sigma_{t-1}^2.
\]
$\gamma > 0$ implies that bad news raises future volatility more
than good news of the same magnitude.
\end{rembox}

% =========================================================
\section*{Summary}
\addcontentsline{toc}{section}{Summary}

\noindent\textbf{Key takeaways.}
\begin{itemize}
  \item Financial returns are uncorrelated but their squares are
        autocorrelated -- volatility comes in clusters.
  \item ARCH$(q)$: $\sigma_t^2$ is a linear function of past
        squared shocks. GARCH$(p, q)$ adds memory through past
        conditional variances.
  \item GARCH$(1, 1)$ is stationary iff $\alpha_1 + \beta_1 < 1$,
        with unconditional variance
        $\alpha_0 / (1 - \alpha_1 - \beta_1)$.
  \item Engle's LM test detects ARCH effects in residuals;
        standardised residuals validate the fit.
  \item Conditional volatility forecasts mean-revert to the
        unconditional variance at rate $\alpha_1 + \beta_1$;
        Value-at-Risk uses these forecasts directly.
  \item Extensions: EGARCH and GJR-GARCH for asymmetric (leverage)
        effects.
\end{itemize}


% =========================================================
% Backmatter: indicative bibliography
% =========================================================
\backmatter

\chapter*{Indicative references}
\addcontentsline{toc}{chapter}{Indicative references}

\noindent\textbf{Foundations and decomposition.}
\begin{itemize}
  \item P.~J.~Brockwell, R.~A.~Davis, \emph{Introduction to Time
        Series and Forecasting}, Springer, 3rd edition, 2016.
  \item C.~Chatfield, \emph{The Analysis of Time Series: An
        Introduction}, Chapman \& Hall/CRC, 7th edition, 2019.
  \item R.~J.~Hyndman, G.~Athanasopoulos, \emph{Forecasting:
        Principles and Practice}, OTexts, 3rd edition, 2021
        (open-access companion: \texttt{otexts.com/fpp3/}).
\end{itemize}

\noindent\textbf{Stationarity, ARMA and forecasting.}
\begin{itemize}
  \item G.~E.~P.~Box, G.~M.~Jenkins, G.~C.~Reinsel, G.~M.~Ljung,
        \emph{Time Series Analysis: Forecasting and Control},
        Wiley, 5th edition, 2015.
  \item P.~J.~Brockwell, R.~A.~Davis, \emph{Time Series: Theory
        and Methods}, Springer, 2nd edition, 2009.
  \item J.~D.~Hamilton, \emph{Time Series Analysis}, Princeton
        University Press, 1994.
  \item R.~H.~Shumway, D.~S.~Stoffer, \emph{Time Series Analysis
        and Its Applications: With R Examples}, Springer, 4th
        edition, 2017.
\end{itemize}

\noindent\textbf{Unit roots and non-stationarity.}
\begin{itemize}
  \item D.~A.~Dickey, W.~A.~Fuller, \emph{Distribution of the
        Estimators for Autoregressive Time Series with a Unit
        Root}, Journal of the American Statistical Association,
        1979.
  \item P.~C.~B.~Phillips, P.~Perron, \emph{Testing for a Unit
        Root in Time Series Regression}, Biometrika, 1988.
  \item D.~Kwiatkowski, P.~C.~B.~Phillips, P.~Schmidt, Y.~Shin,
        \emph{Testing the Null Hypothesis of Stationarity Against
        the Alternative of a Unit Root}, Journal of Econometrics,
        1992.
  \item C.~W.~J.~Granger, P.~Newbold, \emph{Spurious Regressions
        in Econometrics}, Journal of Econometrics, 1974.
\end{itemize}

\noindent\textbf{ARCH, GARCH and volatility.}
\begin{itemize}
  \item R.~F.~Engle, \emph{Autoregressive Conditional
        Heteroscedasticity with Estimates of the Variance of
        United Kingdom Inflation}, Econometrica, 1982.
  \item T.~Bollerslev, \emph{Generalized Autoregressive
        Conditional Heteroskedasticity}, Journal of Econometrics,
        1986.
  \item L.~R.~Glosten, R.~Jagannathan, D.~E.~Runkle, \emph{On the
        Relation between the Expected Value and the Volatility of
        the Nominal Excess Return on Stocks}, Journal of
        Finance, 1993.
  \item D.~B.~Nelson, \emph{Conditional Heteroskedasticity in
        Asset Returns: A New Approach}, Econometrica, 1991.
  \item B.~Mandelbrot, \emph{The Variation of Certain Speculative
        Prices}, Journal of Business, 1963.
  \item P.~H.~Kupiec, \emph{Techniques for Verifying the
        Accuracy of Risk Measurement Models}, Journal of
        Derivatives, 1995.
\end{itemize}

\noindent\textbf{Software documentation.}
\begin{itemize}
  \item S.~Seabold, J.~Perktold, \emph{statsmodels:
        Econometric and Statistical Modeling with Python}, Proc.\
        9th Python in Science Conf., 2010
        (\texttt{statsmodels.org}).
  \item K.~Sheppard, \emph{ARCH Toolbox for Python}
        (\texttt{arch} package, \texttt{bashtage.github.io/arch}).
\end{itemize}

\end{document}
